% !TeX root = ../../Book1.tex
% 纲要标题：## 第9章　曲面联络、测地线与 Gauss–Bonnet
% 纲要位置：几何学计划纲要.md，第 461 行。
% 正文状态：初稿；尚未正式审核。

% 纲要内容开始
% 
% > 本章目标与证明范围：从切向投影建立联络，以平行移动和转角说明曲率的整体意义。二维 Gauss–Bonnet 包括边界和角点的完整证明，不借用后置高维示性类理论。
% 
% 纲要内容结束

\chapter{曲面联络、测地线与 Gauss–Bonnet}
\label{g1:ch:09}
\BookChapterNumber{g1:ch:09}

\begin{chapterabstract}
	沿曲面移动切向量，需要比较不同切平面中的方向。协变导数把这种比较写成平行移动方程，联络一形式则记录移动标架自身的转动。本章建立测地线的初值理论、长度与能量的第一变分以及径向测地线的局部最短性，再把边界切向转角分解为测地曲率与联络积分。局部曲率公式经有限圆盘分解、公共边抵消和顶点计数，成为含边界及角点的 Gauss--Bonnet 定理。最后讨论曲率符号、亏格、向量场零点指数和改变度量时的总曲率不变性。
\end{chapterabstract}

\begin{learninggoals}
	\begin{itemize}
		\item 从切向投影和度量相容性理解沿曲线的协变导数，计算平行移动及绕环路后的旋转。
		\item 证明测地线初值解的存在唯一性与光滑依赖，区分长度、能量、参数选择和端点条件。
		\item 用第一变分与 Gauss 引理说明测地线为何在足够短的范围内最短，并明确局部结论的范围。
		\item 固定联络、曲率、边界方向和外角的符号，完整证明局部及整体 Gauss--Bonnet 公式。
		\item 将总曲率与 Euler 示性数、孤立零点指数联系起来，计算保角改变度量后的内部与边界补偿。
	\end{itemize}
\end{learninggoals}

在球面上沿三个大圆弧走出一个三角形，即使每条边都没有在曲面内主动转弯，最后的方向仍可能与出发时不同。这个现象同时提出了两个问题：沿曲面“保持方向”应如何定义？一条闭合边界的转角又怎样记录它围住的几何？第八章已把曲率写成度量的局部表达式，本章要把它与沿路比较方向、围成区域以及整张曲面的拓扑联系起来。

\ProseParagraphBreak

我们采用 \(\omega(V)=\langle\nabla_Ve_1,e_2\rangle\) 的联络约定，所以曲率关系是 \(\mathrm{d}\omega=-K\,\mathrm{d}A\)。边界按区域在左的方向行进，内角 \(\beta\) 对应外角 \(\alpha=\pi-\beta\)。先在一张坐标图中的圆盘上，借平面转角定理与 Green 公式证明曲率积分公式；随后真正构造分片光滑的有限圆盘分解，把公共边和新增角逐项消去。这样得到的 Gauss--Bonnet 定理包括非定向曲面、光滑边界和有限角点。\footnote{Bonnet，Pierre Ossian Bonnet（1819--1892），法国数学家，研究曲面微分几何、测地曲率与曲率积分公式。}

\ProseParagraphBreak

阅读所需的主要旧结论是第五章圆周角度提升与绕数、第六章曲面分类和 Euler 计数、第七章平面转角定理，以及第八章的度量、协变导数和曲率。使用时会回引具体结果并重述关键条件。常微分方程的光滑依赖、测地线的变分论证与二维 Green 公式所需版本在本章补出；一般流形的微分形式理论和高维示性类理论不是前提。相关教材阅读位置为 Pressley 的平行移动部分\cite[第7.4节]{Pressley2010}、测地线部分\cite[第9.1--9.5节]{Pressley2010}和 Gauss--Bonnet、环路平行移动与向量场零点部分\cite[第13.1--13.4, 13.6--13.7节]{Pressley2010}。

% !TeX root = ../../Book1.tex
% 纲要标题：### 9.1 协变微分和曲面上的移动标架
% 纲要位置：几何学计划纲要.md，第 465 行。
% 正文状态：初稿；尚未正式审核。

% 纲要内容开始
% 
% * 切向投影与度量相容性
% * 沿曲线的平行移动
% * 正交标架的联络一形式
% * 曲率二形式与局部 holonomy
% 
% 纲要内容结束

\section{协变微分和曲面上的移动标架}
\label{g1:sec:0901}

在平面内移动一个箭头，可以要求它在固定直角坐标中的分量不变。曲面上的箭头却必须始终留在当前位置的切平面内；切平面一变，环境空间中的常向量通常就不再合适。第八章的切向投影给出了另一种要求：允许箭头为适应切平面而向法向改变，但不让它在切平面内自行转动。沿路保持这种状态，就是本节要建立的平行移动。

\subsection{切向投影与度量相容性}
\label{g1:subsec:0901:01}

先说明本章的曲面范围。第六章的拓扑曲面默认 Hausdorff、第二可数。给这样的曲面选定坐标图，要求重叠处的坐标变换光滑且具有光滑逆映射，就得到本章所说的\term[guanghuaqumian]{光滑曲面}[smooth surface]。坐标中的切向量按坐标变换的微分换基；光滑正定矩阵 \((g_{ij})\) 若按第八章\eqref{g1:eq:surface-metric-change}在重叠处变换，就在各切平面上给出同一个内积 \(g\)。这样的 \((M,g)\) 称为\term[limanduliangqumian]{黎曼度量曲面}[Riemannian surface]。这里的英文指带黎曼度量的实二维曲面，不是作为复一维对象的 Riemann surface。

\ProseParagraphBreak

正则曲面在 \(\mathbb R^3\) 中的诱导度量是主要例子。不过，平面上按平移识别得到的环面也有相容的欧氏度量；讨论它的内在几何时，不必先把这张带度量的曲面放进三维空间。有光滑边界时用半平面坐标，要求度量系数光滑延拓到坐标域的邻域。后面还会考虑曲面内具有有限角点的区域，边界各光滑段正则，角点内角在 \((0,2\pi)\) 内，因而不包括尖点。

\ProseParagraphBreak

以下指标 \(i,j,k\) 取 \(1,2\)，重复指标求和。回顾\cref{g1:prop:surface-connection-properties}：给定度量后，唯一同时满足度量相容和无挠性的协变导数为
\[
	\nabla_VW=V^i\bigl(\partial_iW^k+\Gamma^k_{ij}W^j\bigr)\partial_k,
	\qquad
	\Gamma^k_{ij}=\frac12g^{k\ell}(\partial_i g_{j\ell}+\partial_j g_{i\ell}-\partial_\ell g_{ij}).
\]
这里 \((g^{ij})=(g_{ij})^{-1}\)。第八章已经证明这套公式在换坐标后给出同一个运算；对于诱导度量，它就是 \(\nabla_VW=(\mathrm{d}W(V))^\top\)。因此本章下面的坐标计算也适用于只给定度量的曲面。

\begin{definition}
	\label{g1:def:surface-along-derivative}
	设 \(\gamma:I\rightarrow M\) 是光滑曲线，\(W(t)\in \mathrm{T}_{\gamma(t)}M\) 是沿曲线的光滑切向量场。在坐标中写 \(\gamma(t)=(u^1(t),u^2(t))\)、\(W(t)=W^k(t)\partial_k|_{\gamma(t)}\)，定义\term[yanquxianxiebiandaoshu]{沿曲线的协变导数}[covariant derivative along a curve]为
	\begin{equation}
		\label{g1:eq:along-covariant-derivative}
		\mathrm{D}_tW\coloneqq\left(\frac{\mathrm{d}W^k}{\mathrm{d}t}+\Gamma^k_{ij}(\gamma(t))\dot u^iW^j\right)\partial_k.
	\end{equation}
	若曲面已浸入 \(\mathbb R^3\)，此式就是 \(\mathrm{D}_tW=(\mathrm{d}W/\mathrm{d}t)^\top\)。\symidx{\(\mathrm{D}_tW\)：沿曲线的协变导数}
\end{definition}

沿曲线的向量场只给每个参数值配置一个向量。即使曲线经过同一点两次，这两次的向量也可以不同，因此不必把 \(W\) 延拓成曲面邻域内的单值向量场。换坐标时，对 \(W\) 的分量应用链式法则，再代入 Christoffel 系数的变换式，二阶坐标导数项相消，故 \(\mathrm{D}_tW\) 仍是同一个切向量。

\begin{proposition}
	\label{g1:prop:along-derivative-properties}
	在光滑度量曲面上，沿光滑曲线的协变导数满足乘法法则及
	\[
		\frac{\mathrm{d}}{\mathrm{d}t}g(V,W)=g(\mathrm{D}_tV,W)+g(V,\mathrm{D}_tW).
	\]
	对光滑双参数映射 \(F(s,t)\)，有 \(\mathrm{D}_s\partial_tF=\mathrm{D}_t\partial_sF\)。若曲面已定向，以 \(J\) 表示每个切平面内的正向直角旋转，则 \(\mathrm{D}_t(JW)=J(\mathrm{D}_tW)\)。
\end{proposition}

\begin{proof}
	乘法法则直接来自\eqref{g1:eq:along-covariant-derivative}。内积求导时，用第八章已证的 \(\partial_i g_{j\ell}=\Gamma^k_{ij}g_{k\ell}+\Gamma^k_{i\ell}g_{jk}\) 将度量系数的导数分到两个向量上，即得第一式。混合导数的第 \(k\) 个分量分别为
	\[
		\partial_s\partial_tF^k+\Gamma^k_{ij}(F)\partial_sF^i\partial_tF^j,
		\qquad
		\partial_t\partial_sF^k+\Gamma^k_{ij}(F)\partial_tF^i\partial_sF^j.
	\]
	偏导可交换，且 \(\Gamma^k_{ij}=\Gamma^k_{ji}\)，所以两式相等；\(F\) 不必是浸入。

	最后局部选正向正交单位标架 \((e_1,e_2)\)。由内积求导，存在标量 \(a(t)\)，使 \(\mathrm{D}_te_1=ae_2\)、\(\mathrm{D}_te_2=-ae_1\)。将 \(W=xe_1+ye_2\)、\(JW=-ye_1+xe_2\) 分别求导，便得 \(\mathrm{D}_t(JW)=J(\mathrm{D}_tW)\)。
\end{proof}

这里的 \(J\) 仅用度量与定向定义，满足 \(J^2=-\mathrm{id}\)；在取单位法向 \(n\) 的空间曲面上，\(JW=n\times W\)。以后边界的符号都以这个正向直角旋转为准。

\subsection{沿曲线的平行移动}
\label{g1:subsec:0901:02}

\begin{definition}
	沿分片光滑曲线 \(\gamma\) 的连续切向量场 \(W\)，若在每个光滑段上满足 \(\mathrm{D}_tW=0\)，就称为沿 \(\gamma\) 的\term[pingxingxiangliangchang]{平行向量场}[parallel vector field]。从给定初向量出发求出终向量的操作称为\term[pingxingyidong]{平行移动}[parallel transport]，记为 \(P_\gamma\)。\symidx{\(P_\gamma\)：沿曲线的平行移动}
\end{definition}

\begin{theorem}
	\label{g1:thm:surface-parallel-transport}
	设 \((M,g)\) 是光滑度量曲面，\(\gamma:[a,b]\rightarrow M\) 分片光滑。给定 \(w\in \mathrm{T}_{\gamma(a)}M\)，存在唯一沿 \(\gamma\) 的平行场，初值为 \(w\)。所得 \(P_\gamma:\mathrm{T}_{\gamma(a)}M\rightarrow \mathrm{T}_{\gamma(b)}M\) 是线性等距同构；在定向曲面上还保持定向。它在保向重参数化下不变，并满足
	\[
		P_{\gamma_1*\gamma_2}=P_{\gamma_2}\circ P_{\gamma_1},
		\qquad P_{\bar\gamma}=P_\gamma^{-1}.
	\]
	这里沿用第五章的次序，\(\gamma_1*\gamma_2\) 先走 \(\gamma_1\) 再走 \(\gamma_2\)，\(\bar\gamma\) 表示反向道路。
\end{theorem}

\begin{proof}
	在一张坐标图中，平行条件是线性方程
	\[
		\dot W^k=-\Gamma^k_{ij}(\gamma(t))\dot u^iW^j.
	\]
	每个光滑闭段上的系数连续，故\cref{g1:lem:linear-frame-ode}给出指定初值的唯一解。把曲线的坐标邻域拉回参数区间，利用紧区间的有限细分，使每段落入一张图。在接点把所得几何向量换成下一张图的分量作为初值；坐标无关性与唯一性保证相邻解相容，也保证改变细分不会改变结果。

	方程线性，所以解依初值线性。对两个平行场，由上命题得 \(\frac{\mathrm{d}}{\mathrm{d}t}g(V,W)=0\)，于是长度与夹角保持。反向参数化时 \(\mathrm{D}_\tau W=t'(\tau)\mathrm{D}_tW\)，故平行性不变；唯一性给出反向和拼接公式，也给出可逆性。在定向情形，把一个正向单位标架平行移动，两向量始终正交且单位长；其方向符号连续而不可能经过零，故一直保持正向。
\end{proof}

\begin{example}
	\label{g1:exm:parallel-polar-frame}
	在欧氏平面沿圆 \(\gamma(t)=(r\cos t,r\sin t)\) 运动，取 \(e_r=(\cos t,\sin t)\)、\(e_\theta=(-\sin t,\cos t)\)。若 \(W=a e_r+b e_\theta\)，则
	\[
		\mathrm{D}_tW=(a'-b)e_r+(b'+a)e_\theta.
	\]
	由 \(W(0)=(1,0)\) 出发的平行场满足 \(a=\cos t\)、\(b=-\sin t\)，所以 \(W(t)=(1,0)\) 在环境平面内确实不变。反而是分量固定的 \(e_r\) 本身不平行：它一直朝向圆外，因而在转动。分量变化并不等于向量发生了内在转动，必须同时计算标架的变化。
\end{example}

\subsection{正交标架的联络一形式}
\label{g1:subsec:0901:03}

在一张定向坐标图中，先对 \(\partial_1,\partial_2\) 作 Gram--Schmidt 正交化，得到光滑正向正交标架 \((e_1,e_2)\)。内积为常数意味着 \(\nabla_Ve_1\) 沿 \(e_2\)，而 \(\nabla_Ve_2\) 沿 \(-e_1\)。因此所有方向上的标架转动，只需用一个对 \(V\) 线性的数来记录。

\begin{definition}
	\label{g1:def:surface-connection-form}
	对定向光滑度量曲面上的局部正向正交标架，定义\term[lianluoyixingshi]{联络一形式}[connection one-form]为 \(\omega(V)=g(\nabla_Ve_1,e_2)\)。于是
	\begin{equation}
		\label{g1:eq:surface-frame-connection}
		\nabla_Ve_1=\omega(V)e_2,\qquad \nabla_Ve_2=-\omega(V)e_1.
	\end{equation}
\end{definition}

这里“一形式”表示在每点给定一个随点光滑变化的线性函数。在坐标中它就是 \(A\,\mathrm{d}u+B\,\mathrm{d}v\)，其中 \(\mathrm{d}u(V)=V^u\)。两一形式的外积定义为 \((\alpha\wedge\beta)(V,W)=\alpha(V)\beta(W)-\alpha(W)\beta(V)\)，因而 \(\mathrm{d}u\wedge\mathrm{d}v\) 记录两个向量的有向坐标面积。令 \(\vartheta^i(V)=g(V,e_i)\)，则
\[
	\mathrm{d}A=\vartheta^1\wedge\vartheta^2=\sqrt{\det(g_{ij})}\,\mathrm{d}u\wedge\mathrm{d}v
\]
是正向面积形式。计算非负面积时取其对应的面积元；反向坐标使用 Jacobian 行列式的绝对值。\symidx{\(\mathrm{d}A\)：曲面的面积元}

对一形式定义
\[
	\mathrm{d}(A\,\mathrm{d}u+B\,\mathrm{d}v)
	=(\partial_uB-\partial_vA)\,\mathrm{d}u\wedge\mathrm{d}v.
\]
对函数 \(f\) 则有 \(\mathrm{d}f=f_u\,\mathrm{d}u+f_v\,\mathrm{d}v\)，混合偏导相等给出 \(\mathrm{d}(\mathrm{d}f)=0\)。这些是本章使用的\term[waiweifen]{外微分}[exterior derivative]。换坐标时应用链式法则，二阶坐标导数项两两相消；更具几何意味的同一公式是
\[
	\mathrm{d}\alpha(V,W)=V(\alpha(W))-W(\alpha(V))-\alpha([V,W]).
\]
将 \(V,W\) 的坐标分量展开可验证此式，也直接看出其与坐标无关。我们在这里需要的只是函数、一形式和面积形式，不以掌握高维微分形式为前提。

\begin{proposition}
	\label{g1:prop:connection-frame-change}
	若另一正向正交标架在邻域内写成
	\[
		\widetilde e_1=\cos\varphi\,e_1+\sin\varphi\,e_2,\qquad
		\widetilde e_2=-\sin\varphi\,e_1+\cos\varphi\,e_2,
	\]
	则相应联络形式满足 \(\widetilde\omega=\omega+\mathrm{d}\varphi\)。沿曲线写单位向量 \(W=\cos\theta\,e_1+\sin\theta\,e_2\)，有
	\begin{equation}
		\label{g1:eq:parallel-angle}
		\mathrm{D}_tW=(\theta'+\omega(\dot\gamma))JW.
	\end{equation}
	特别地，平行条件为 \(\theta'=-\omega(\dot\gamma)\)。
\end{proposition}

\begin{proof}
	将两个标架变换式求协变导数，用\eqref{g1:eq:surface-frame-connection}合并正弦、余弦项，得到 \(\nabla_V\widetilde e_1=(V\varphi+\omega(V))\widetilde e_2\)，即第一式。对 \(W\) 作相同展开，得到第二式。沿参数区间的连续角度由\cref{g1:lem:circle-angle-lift}选取；只要 \(W\) 光滑，局部反三角函数又说明提升角度光滑。
\end{proof}

联络形式依赖标架；标架多转了 \(\varphi\)，所记录的转动便多出 \(\mathrm{d}\varphi\)。但 \(\mathrm{d}\widetilde\omega=\mathrm{d}\omega\)。在环带上换架角度未必有全域单值的实函数表示，上述关系仍在每个小邻域成立，外微分不依赖标架的结论因而仍然成立。

\subsection{曲率二形式与局部 holonomy}
\label{g1:subsec:0901:04}

\begin{proposition}
	\label{g1:prop:connection-curvature-form}
	在定向光滑度量曲面上，任意局部正向正交标架的联络形式均满足
	\begin{equation}
		\label{g1:eq:surface-curvature-form}
		\mathrm{d}\omega=-K\,\mathrm{d}A.
	\end{equation}
	这里 \(K\) 是第八章定义的内在 Gauss 曲率，曲率张量约定为 \(R(V,W)=\nabla_V\nabla_W-\nabla_W\nabla_V-\nabla_{[V,W]}\)。
\end{proposition}

\begin{proof}
	对 \(e_1\) 应用曲率定义与\eqref{g1:eq:surface-frame-connection}，两次求导中沿 \(e_1\) 的乘积项抵消，得到
	\[
		R(V,W)e_1=\bigl(V\omega(W)-W\omega(V)-\omega([V,W])\bigr)e_2
		=\mathrm{d}\omega(V,W)e_2.
	\]
	另一方面，\cref{g1:prop:surface-curvature-two-dim}给出 \(R(e_1,e_2)e_1=-Ke_2\)。所以 \(\mathrm{d}\omega(e_1,e_2)=-K\)。二维交替双线性形式由在一组基上的这个值决定，而 \(\mathrm{d}A(e_1,e_2)=1\)，故得结论。
\end{proof}

联络的曲率矩阵是 \(\begin{psmallmatrix}0&-\mathrm{d}\omega\\\mathrm{d}\omega&0\end{psmallmatrix}\)；它记录了先后两个方向求导的差。负号 \(\mathrm{d}\omega=-K\,\mathrm{d}A\) 同时来自本书的曲率次序和 \(\omega=g(\nabla e_1,e_2)\) 的选择。改变其中任何一项约定时，都必须重新核对这个符号。

\begin{lemma}[Green 公式]
	\label{g1:lem:green-formula}
	设 \(D\) 是一张正向坐标图内的紧区域，边界由有限条分片 \(C^2\) 正则简单闭曲线组成，彼此不交；角点内角在 \((0,2\pi)\) 内。若一形式 \(\alpha\) 在 \(D\) 的邻域上为 \(C^1\)，则
	\[
		\int_{\partial D}\alpha=\int_D\mathrm{d}\alpha.
	\]
	边界方向取区域在左的方向，故有孔时，外边界逆时针，孔的边界顺时针。
\end{lemma}

这一公式称为 Green 公式，也就是二维的 Stokes 公式。\footnote{Green，格林，George Green（1793--1841），英国数学家，研究位势理论，提出 Green 函数及相关积分公式。Stokes，斯托克斯，George Gabriel Stokes（1819--1903），生于爱尔兰的英国数学家、物理学家，研究流体力学与光学。}下面用平面微积分证明本章所需的版本。

\begin{proof}
	先设边界光滑。在边界附近选坐标 \((u,t)\)，使区域为 \(t\ge0\)、边界为 \(t=0\)，并保持定向；这由边界正则性及逆函数定理得到。例如边界原为 \(v=f(u)\)，取 \(t=v-f(u)\)。若 \(\alpha=P\,\mathrm{d}u+Q\,\mathrm{d}t\) 的支撑紧含于此坐标邻域，则积分基本定理与累次积分给出
	\[
		\int_{t\ge0}(\partial_uQ-\partial_tP)\,\mathrm{d}u\,\mathrm{d}t
		=\int P(u,0)\,\mathrm{d}u.
	\]
	第一项沿 \(u\) 积分为零，第二项只留下 \(t=0\) 的值；右边正是把区域放在左侧的边界积分。若支撑在内部，同样在整个平面上积分，两个偏导积分都为零。

	把 \(D\) 用有限个内部邻域和这样的边界邻域覆盖，再用支撑在其中的光滑函数 \(\rho_j\) 分解 \(1\)。这里不需调用曲面上的一般分解定理：在坐标平面内取有限个更小圆盘覆盖紧集，以 \(\exp(-1/(1-|x|^2))\) 在单位圆盘内、在外取零并经平移缩放得到非负光滑函数，最后除以它们在紧集附近处处为正的和即可。对 \(\rho_j\alpha\) 应用局部计算并相加，由 \(\sum\rho_j=1\)、\(\sum\mathrm{d}\rho_j=0\) 得到原式。换坐标的积分公式保证所有局部积分相容。

	有角点时，先在每个角点附近用一条长度为 \(O(\varepsilon)\) 的光滑小弧切去角点，并把新接点圆滑化。两侧边界有非零夹角，可在其切线张成的扇形内完成此构造；所删区域面积为 \(O(\varepsilon^2)\)，所改边界的总长度为 \(O(\varepsilon)\)。\(\alpha\) 及 \(\mathrm{d}\alpha\) 在紧邻域上有界，所以新旧边界积分之差趋于零，面积积分之差也趋于零。对光滑截角区域的公式取极限，便得分片光滑情形；这一论证也适用于多个边界分支。
\end{proof}

对任意紧曲面区域，面积积分也可以先作有限坐标分解再相加。在曲面坐标中取上述非负光滑截断函数，支撑紧含于坐标邻域后延为零，再除以有限和，得到坐标覆盖的有限单位分解。以局部表达计算 \(\int\rho_j f\,\mathrm{d}A\)；换坐标的 Jacobian 公式及 \(\sum\rho_j=1\) 保证结果与分解无关。面积元在不可定向曲面上仍有意义，因为它使用绝对面积；带方向的二形式积分则需要定向。

\begin{proposition}
	\label{g1:prop:surface-local-holonomy}
	设 \(D\) 是上述坐标图内由简单分片光滑闭曲线围成的圆盘区域，\(\gamma=\partial D\) 取正向。沿 \(\gamma\) 的平行移动回到起点后是一个平面旋转，其旋转角 \(\beta\) 满足
	\begin{equation}
		\label{g1:eq:surface-holonomy-integral}
		\beta\equiv\int_D K\,\mathrm{d}A\pmod{2\pi}.
	\end{equation}
\end{proposition}

\begin{proof}
	在整张坐标图上将坐标基正交化，得到单值标架。绕行一周，标架回到原值。平行单位向量相对标架的连续角度由\eqref{g1:eq:parallel-angle}满足 \(\theta(b)-\theta(a)=-\int_{\partial D}\omega\)。再用 Green 公式和\eqref{g1:eq:surface-curvature-form}，得到角度增量 \(\int_DK\,\mathrm{d}A\)。终点向量只记录该数模 \(2\pi\) 的值，这便是所述旋转角。
\end{proof}

绕闭路返回同一切平面后所得的变换也称为该闭路的 holonomy。\foreignidx{holonomy}它通常不只取决于闭路的同伦类：同一个坐标圆盘内的闭路都可缩，但围起不同曲率积分的闭路可能给出不同旋转。对缩到点 \(p\) 的小圆盘，选取趋于零的角度代表后，有
\[
	\frac{\beta}{\operatorname{Area}(D)}
	=\frac{1}{\operatorname{Area}(D)}\int_DK\,\mathrm{d}A\longrightarrow K(p),
\]
因为曲率的连续性控制了平均值与 \(K(p)\) 的差。这给出了曲率的一种直接测量方式。

\begin{example}
	\label{g1:exm:spherical-latitude-holonomy}
	半径 \(R\) 的球面用余纬 \(\theta\in(0,\pi)\)、经度 \(\varphi\) 表示，度量为 \(R^2\mathrm{d}\theta^2+R^2\sin^2\theta\,\mathrm{d}\varphi^2\)。取 \(e_1=R^{-1}\partial_\theta\)、\(e_2=(R\sin\theta)^{-1}\partial_\varphi\)，由 Christoffel 公式求得
	\[
		\omega=\cos\theta\,\mathrm{d}\varphi,\qquad
		\mathrm{d}\omega=-\sin\theta\,\mathrm{d}\theta\wedge\mathrm{d}\varphi=-R^{-2}\mathrm{d}A.
	\]
	沿余纬 \(\theta_0\) 处的纬线增加经度一周，平行向量相对这个周期标架的角度变化为 \(-2\pi\cos\theta_0\)。而北球冠的曲率积分是 \(2\pi(1-\cos\theta_0)\)，两者相差 \(2\pi\)，表示同一个旋转。这里的经纬标架在北极无定义，不能将它的 \(\omega\) 直接用于整个球冠的 Green 公式；换用在球冠上光滑的坐标标架，角度代表就与球冠积分一致。赤道绕行的旋转为恒等，但北半球的曲率积分为 \(2\pi\)，这也说明模 \(2\pi\) 不能随意删去。
\end{example}

% !TeX root = ../../Book1.tex
% 纲要标题：### 9.2 测地线与变分
% 纲要位置：几何学计划纲要.md，第 472 行。
% 正文状态：初稿；尚未正式审核。

% 纲要内容开始
% 
% * 测地线方程的存在唯一性
% * 弧长、能量与端点条件
% * 测地曲率及第一变分
% * 测地坐标和径向局部极小性
% 
% 纲要内容结束

\section{测地线与变分}
\label{g1:sec:0902}

在平面上，匀速直线同时具有两种特征：速度沿运动保持不变，较短的直线段又是两端之间的最短路。曲面上不宜直接要求空间中的速度不变，因为这个向量可能立即离开切平面。上一节的平行移动给出了替代条件：让速度沿曲线平行。它定义的是一个微分方程；这一方程与长度极小之间的关系还需要证明，而且最短性首先只在短程成立。

\ProseParagraphBreak

以下讨论的曲面 \(M\) 都带光滑正定度量：可以取正则曲面及其诱导度量，也可以取上一节所说的二维光滑度量曲面。局部计算取坐标 \((u^1,u^2)\)，记 \(g_{ij}=\langle\partial_i,\partial_j\rangle\)。讨论初值与邻域时，所取点均为曲面的内点；边界会在变分的端点条件中单独出现。

\subsection{测地线方程的存在唯一性}
\label{g1:subsec:0902:01}

\begin{definition}
	\label{g1:def:geodesic}
	光滑度量曲面 \(M\) 上的光滑曲线 \(\gamma:I\to M\)，若满足
	\[
		\mathrm{D}_t\dot\gamma=0,
	\]
	就称为\term[cedixian]{测地线}[geodesic]。这里 \(\mathrm{D}_t\) 是\cref{g1:def:surface-along-derivative}定义的沿曲线协变导数，参数 \(t\) 称为仿射参数。常值曲线也包括在内。
\end{definition}

由\cref{g1:prop:along-derivative-properties}的度量相容性，测地线满足
\[
	\frac{\mathrm{d}}{\mathrm{d}t}|\dot\gamma|^2=2\langle \mathrm{D}_t\dot\gamma,\dot\gamma\rangle=0.
\]
所以每条非常值测地线都以恒定的正速度行进，改用弧长后仍是测地线。若 \(\gamma\) 是非常值测地线，重参数曲线 \(\eta(s)=\gamma(\phi(s))\) 满足
\[
	\mathrm{D}_s\dot\eta=\phi''(s)\dot\gamma(\phi(s)),
\]
因而保持测地线方程的正则重参数化恰为 \(\phi(s)=as+b\)、\(a\ne0\)。同一条轨迹以忽快忽慢的速度走过，便可能不再满足上述带参数的定义。

\ProseParagraphBreak

把沿曲线求导的坐标式代入定义，测地线方程成为
\begin{equation}
	\label{g1:eq:geodesic-equation}
	\ddot u^k+\Gamma^k_{ij}(u)\dot u^i\dot u^j=0,\qquad k=1,2.
\end{equation}
这里对 \(i,j\) 求和。\(\Gamma^k_{ij}\) 由\eqref{g1:eq:christoffel-metric}计算，故这一方程只依赖于度量。在嵌入曲面上，它也表示空间加速度全部沿法向；一般并不表示空间加速度为零。

\begin{lemma}[初值的光滑依赖]
	\label{g1:lem:smooth-ode-flow}
	设 \(F:\Omega\to\mathbb R^m\) 是开集 \(\Omega\subset\mathbb R^m\) 上的光滑映射，\(z_0\in\Omega\)。则存在 \(a>0\) 及 \(z_0\) 的邻域 \(B\)，使
	\[
		\partial_t\Phi(t,z)=F(\Phi(t,z)),\qquad \Phi(0,z)=z
	\]
	对 \((t,z)\in(-a,a)\times B\) 有唯一局部解，且 \(\Phi\) 对时间与初值联合光滑。若一条解已经存在于紧时间区间，则该区间上的解也对其附近的初值光滑依赖。
\end{lemma}
\begin{proof}
	\cref{g1:lem:surface-field-integral-curve}用积分迭代证明了二维向量场的局部存在唯一性。该证明只使用 Euclid 范数与导数界，在 \(\mathbb R^m\) 中同样成立。具体地，取 \(\overline B(z_0,2r)\subset\Omega\)，在此球上以 \(M,L\) 分别控制 \(|F|,\|\mathrm{D}F\|\)。令初值 \(z\in B(z_0,r)\)，并取 \(aM<r\)、\(q=aL<1\)，则
	\[
		(\mathcal T_zx)(t)=z+\int_0^tF(x(s))\,\mathrm{d}s
	\]
	把取值于大闭球的连续路径映入自身，且一致距离至多缩小为 \(q\) 倍。相邻迭代的差构成收敛几何级数，极限就是唯一解。比较两个初值得到
	\[
		\|\Phi(\cdot,z+h)-\Phi(\cdot,z)\|_\infty\le\frac{|h|}{1-q}.
	\]
	这也先证明了初值依赖的连续性。

	为证明可微，令矩阵 \(A(t,z)\) 解积分方程
	\[
		A(t,z)=I+\int_0^t \mathrm{D}F(\Phi(s,z))A(s,z)\,\mathrm{d}s.
	\]
	同一个缩短估计给出其唯一连续解，并由比较积分方程得到它对 \(z\) 连续。记 \(\delta_h=\Phi(\cdot,z+h)-\Phi(\cdot,z)\)。\(\mathrm{D}F\) 在较小闭球上一致连续，Taylor 余项与上面的 Lipschitz 估计遂给出
	\[
		\|\delta_h-A(\cdot,z)h\|_\infty
		\le q\|\delta_h-A(\cdot,z)h\|_\infty+o(|h|).
	\]
	把第一项移到左边，就有 \(\mathrm{D}_z\Phi=A\)。高阶导数可以用相同办法逐阶得到：对初值求 \(k\) 阶导后，未知的最高阶导数 \(B_k\) 总满足
	\[
		B_k(t)=B_k(0)+\int_0^t\bigl(\mathrm{D}F(\Phi(s))B_k(s)+C_k(s)\bigr)\,\mathrm{d}s,
	\]
	其中 \(C_k\) 只由 \(F\) 的至多 \(k\) 阶导数和已经求得的低阶导数组成，\(k\ge2\) 时 \(B_k(0)=0\)。归纳时，对前一阶方程取差商，Taylor 余项趋零，再用 \(1-q>0\) 的同一估计，便证明该方程的解确为下一阶导数而且连续。时间导数则由 \(\partial_t\Phi=F(\Phi)\) 得到，混合导数也随之连续，故 \(\Phi\) 联合光滑。

	最后，已有解在紧时间区间上的像是紧集，可沿它选取有限个上述局部时间步。逐步缩小初值邻域，使各步的终点仍落在下一步的允许邻域；有限次光滑映射的复合给出整个区间上的光滑依赖，唯一性保证相邻步骤相接。
\end{proof}

\begin{proposition}[测地线初值问题]
	\label{g1:prop:geodesic-initial-values}
	设 \(M\) 是光滑度量曲面，\(p\) 为内点，\(v\in \mathrm{T}_pM\)。则存在以 \(0\) 为内点的区间及其上的唯一局部测地线 \(\gamma_v\)，使 \(\gamma_v(0)=p\)、\(\dot\gamma_v(0)=v\)。在共同存在的紧时间区间上，解对起点与初速度光滑依赖。若 \(v\ne0\)，其速度恒为 \(|v|\)；若 \(v=0\)，则为常值曲线。
\end{proposition}
\begin{proof}
	置 \(w^k=\dot u^k\)，将\eqref{g1:eq:geodesic-equation}改写成四维一阶系统
	\[
		\dot u^k=w^k,\qquad \dot w^k=-\Gamma^k_{ij}(u)w^iw^j.
	\]
	右端光滑，前一引理给出全部初值结论。换坐标后所得曲线仍满足几何定义 \(\mathrm{D}_t\dot\gamma=0\)，所以局部唯一性使各坐标中的解相符。速度性质已由度量相容性证明；零初速度时，常值曲线是一个解，唯一性排除了其他解。
\end{proof}

初值问题给定的是“一点及一个速度”，并不保证任意两点间有唯一测地线，也不保证解能够走任意长时间。例如去掉原点的 Euclid 平面中，指向原点的直线在有限时间撞上被删去的点；而球面的南北两极之间有许多不同的半大圆。

\begin{example}[球面的测地线]
	\label{g1:exm:sphere-geodesics}
	半径 \(R\) 的球面取外向单位法向 \(n=\gamma/R\)。若 \(\gamma\) 是速度 \(c>0\) 的测地线，则空间加速度平行于 \(\gamma\)。两次求导 \(\langle\gamma,\gamma\rangle=R^2\)，有 \(\langle\ddot\gamma,\gamma\rangle=-c^2\)，所以
	\[
		\ddot\gamma=-\frac{c^2}{R^2}\gamma,\qquad
		\gamma(t)=p\cos\frac{ct}{R}+\frac{R}{c}v\sin\frac{ct}{R},
	\]
	其中 \(|p|=R\)、\(\langle p,v\rangle=0\)、\(|v|=c\)。这个公式也可直接核验初值与方程，唯一性说明它给出全部非常值测地线。因此球面的测地线轨迹是大圆；绕一整圈仍满足方程，却显然不是起终点之间的最短路。
\end{example}

\subsection{弧长、能量与端点条件}
\label{g1:subsec:0902:02}

比较曲线时，长度天然不在乎走得快慢，微分方程却记录参数。引入能量可以把参数也纳入比较，使随后的求导不出现速度范数的分母。对于固定区间 \([a,b]\)、\(a<b\) 上的分片 \(C^1\) 曲线，定义
\[
	L(\gamma)=\int_a^b|\dot\gamma(t)|\,\mathrm{d}t,\qquad
	E(\gamma)=\frac12\int_a^b|\dot\gamma(t)|^2\,\mathrm{d}t.
\]
这里的 \(E\) 称为曲线的\term[nengliang]{能量}[energy]。\symidx{\(E(\gamma)\)：曲线能量}长度是各坐标段中的度量长度；对于嵌入曲面，\cref{g1:prop:arc-length-param}的长度公式说明它与空间曲线的折线定义一致。

\begin{proposition}[长度与能量]
	\label{g1:prop:length-energy}
	设 \((M,g)\) 为光滑度量曲面，\(a<b\)，\(\gamma:[a,b]\to M\) 为分片 \(C^1\) 曲线。令 \(L(\gamma)=\int_a^b|\dot\gamma|_g\,\mathrm dt\)、\(E(\gamma)=\tfrac12\int_a^b|\dot\gamma|_g^2\,\mathrm dt\)，则
	\[
		L(\gamma)^2\le2(b-a)E(\gamma).
	\]
	等号成立当且仅当速度范数几乎处处为常数；对 \(C^1\) 曲线，这等价于处处等速。若一条曲线在给定端点的全部竞争曲线中长度最小且等速，则它在同一区间上也使能量最小。
\end{proposition}
\begin{proof}
	对函数 \(|\dot\gamma|\) 与常函数 \(1\) 使用 Cauchy--Schwarz 不等式，便得到所述估计。等号条件是两个函数几乎处处成比例。若 \(\gamma_0\) 等速且最短，对任意同端点的竞争曲线 \(\eta\)，则
	\[
		E(\eta)\ge\frac{L(\eta)^2}{2(b-a)}\ge\frac{L(\gamma_0)^2}{2(b-a)}=E(\gamma_0).
	\]
\end{proof}

对正则曲线，也可把其轨迹改为等速行进而不改长度。具体地，令 \(\ell(t)=\int_a^t|\dot\gamma(\tau)|\,\mathrm{d}\tau\)，它严格递增，按\cref{g1:prop:arc-length-param}中正速度积分取逆的构造，再把 \([0,L]\) 线性缩放到 \([a,b]\)，便得到速度 \(L/(b-a)\) 的参数。该构造只使用正的速度函数，对这里的度量速度同样成立。因此在正则曲线的比较中，能量最小同时要求轨迹最短和参数等速；长度最小则不强制后一条件。

\ProseParagraphBreak

为了精确表达“小幅移动一条曲线”，取光滑映射
\[
	F:(-\varepsilon,\varepsilon)\times[a,b]\to M,\qquad F(0,t)=\gamma(t).
\]
称 \(F\) 为 \(\gamma\) 的一个\term[bianfen]{变分}[variation]，称 \(V(t)=\partial_sF(0,t)\) 为变分向量场。固定端点是要求 \(F(s,a)=p\)、\(F(s,b)=q\)，它给出 \(V(a)=V(b)=0\)。若只要求端点分别落在给定的光滑曲线 \(C_0,C_1\) 上，则 \(V(a)\in \mathrm{T}_pC_0\)、\(V(b)\in \mathrm{T}_qC_1\)，这些向量通常不为零。若端点可以在曲面内任意移动，相应的端点向量更没有限制。这三种条件不能混用，下一小节中的边界项会记录它们的差别。

\subsection{测地曲率及第一变分}
\label{g1:subsec:0902:03}

空间曲率把曲线所有弯曲都算在内。对于曲面内部的最短路，真正应受约束的是曲线在切平面内向侧方偏转的程度。以下先取一个曲面定向，记 \(J\) 为切平面内的正向直角旋转；嵌入曲面中，\(JW=n\times W\)。

\begin{definition}
	\label{g1:def:geodesic-curvature}
	设 \((M,g)\) 为定向光滑度量曲面，\(\mathrm D\) 为其无挠相容联络沿曲线的导数，\(J\) 为各切平面的正向直角旋转。设 \(\gamma\) 是 \(M\) 上的 \(C^2\) 正则曲线，\(s\) 为沿行进方向的弧长，\(T=\mathrm{d}\gamma/\mathrm{d}s\)。定义有向\term[cediqulv]{测地曲率}[geodesic curvature]为
	\[
		k_g=\langle \mathrm{D}_sT,JT\rangle.
	\]
	由于 \(|T|=1\)，度量相容性给出 \(\mathrm{D}_sT\perp T\)，故 \(\mathrm{D}_sT=k_gJT\)。\symidx{\(k_g\)：有向测地曲率}
\end{definition}

于是正则曲线的弧长参数化为测地线，当且仅当 \(k_g=0\)。反转曲面定向会使 \(k_g\) 变号；保持曲面定向而反转曲线行进方向也使 \(k_g\) 变号。若 \(\gamma\) 是嵌入曲面上的单位速曲线，\cref{g1:prop:gauss-weingarten}还给出空间加速度的分解
\begin{equation}
	\label{g1:eq:geodesic-normal-curvature}
	\gamma''=k_gJT+\mathrm{II}(T,T)n,
	\qquad \kappa^2=k_g^2+\mathrm{II}(T,T)^2.
\end{equation}
测地线只消去这里的切向项，通常仍有法曲率。例如球面大圆的测地曲率为零，空间曲率却是 \(1/R\)。

\begin{proposition}[第一变分]
	\label{g1:prop:first-variation}
	设 \((M,g)\) 为光滑度量曲面，\(a<b\)、\(\varepsilon>0\)，\(F:(-\varepsilon,\varepsilon)\times[a,b]\to M\) 是光滑变分，\(\gamma(t)=F(0,t)\)、\(V(t)=\partial_sF(0,t)\)。令 \(E(c)=\tfrac12\int_a^b|\dot c|_g^2\,\mathrm dt\)、\(L(c)=\int_a^b|\dot c|_g\,\mathrm dt\)，\(\mathrm D\) 取 \(g\) 的无挠相容联络。则
	\begin{equation}
		\label{g1:eq:energy-first-variation}
		\left.\frac{\mathrm{d}}{\mathrm{d}s}E(F(s,\cdot))\right|_{s=0}
		=\bigl[\langle V,\dot\gamma\rangle\bigr]_a^b-\int_a^b\langle V,\mathrm{D}_t\dot\gamma\rangle\,\mathrm{d}t.
	\end{equation}
	若 \(\gamma\) 正则，令 \(T=\dot\gamma/|\dot\gamma|\)，则
	\begin{equation}
		\label{g1:eq:length-first-variation}
		\left.\frac{\mathrm{d}}{\mathrm{d}s}L(F(s,\cdot))\right|_{s=0}
		=\bigl[\langle V,T\rangle\bigr]_a^b-\int_a^b\langle V,\mathrm{D}_tT\rangle\,\mathrm{d}t.
	\end{equation}
	若 \(M\) 已定向，以 \(J\) 记正向直角旋转，\(k_g=\langle\mathrm D_\ell T,JT\rangle\) 记测地曲率，把最后一项改用弧长 \(\ell\) 积分，即为 \(-\int_\gamma k_g\langle V,JT\rangle\,\mathrm{d}\ell\)。
\end{proposition}
\begin{proof}
	记 \(W=\partial_tF\)。由\cref{g1:prop:along-derivative-properties}，度量相容性及混合协变导数相等分别给出
	\[
		\partial_s\frac12\langle W,W\rangle=\langle \mathrm{D}_sW,W\rangle,
		\qquad \mathrm{D}_s\partial_tF=\mathrm{D}_t\partial_sF.
	\]
	后一个等式的坐标理由是 \(\partial_s\partial_tu^k=\partial_t\partial_su^k\) 及 \(\Gamma^k_{ij}=\Gamma^k_{ji}\)，并不要求曲率为零。对紧区间积分可在积分号内求导，故能量导数是 \(\int_a^b\langle \mathrm{D}_tV,\dot\gamma\rangle\,\mathrm{d}t\)。再由内积的求导公式作分部积分，得到\eqref{g1:eq:energy-first-variation}。

	若 \(\gamma\) 正则，其速度在紧区间上有正的下界，缩小变分参数区间后附近曲线也正则。因此可对速度范数直接求导，得到长度导数
	\[
		\int_a^b\left\langle \mathrm{D}_tV,\frac{\dot\gamma}{|\dot\gamma|}\right\rangle\,\mathrm{d}t.
	\]
	再次分部积分即得\eqref{g1:eq:length-first-variation}。最后 \(\mathrm{D}_tT=|\dot\gamma|\mathrm{D}_\ell T\)，而 \(\mathrm{d}\ell=|\dot\gamma|\mathrm{d}t\)，代入测地曲率定义便得弧长形式。
\end{proof}

所谓驻定，是指对所有允许变分的一阶导数均为零，不是指附近所有曲线都更长。固定端点时，\eqref{g1:eq:energy-first-variation}立即说明测地线使能量驻定。反过来，如果能量对所有固定端点变分驻定，则 \(\mathrm{D}_t\dot\gamma=0\)：若某个内点处该向量非零，取其附近一段落在同一坐标图中的小区间，并选非负光滑函数 \(\chi\)，支撑在此小区间内且在该点为正。令 \(V=\chi \mathrm{D}_t\dot\gamma\)，在坐标中用 \(u_s(t)=u(t)+sV^i(t)\) 实现这个变分，并在区间外保持原曲线，则能量导数为
\[
	-\int_a^b\chi|\mathrm{D}_t\dot\gamma|^2\,\mathrm{d}t<0,
\]
产生矛盾。所需的 \(\chi\) 可由区间内部的函数 \(\exp(-1/(\delta^2-(t-t_0)^2))\) 在区间外补零得到。对长度公式作相同论证，正则曲线的固定端点长度驻定等价于 \(\mathrm{D}_tT=0\)，即轨迹按弧长走时为测地线；它本身未必等速。

\ProseParagraphBreak

若端点分别可在 \(C_0,C_1\) 上移动，先用支撑在内部的变分仍得到测地线条件，再保留边界项：驻定还要求终点速度分别垂直于 \(\mathrm{T}_pC_0,\mathrm{T}_qC_1\)。任意端点切向量都可实现：先在相应端点曲线上取具有该速度的短曲线，再在端点附近的坐标图中用光滑截断函数把这一移动延伸到变分，远处保持不动。因而上述垂直条件也是必要的；反过来，测地线条件和垂直条件使全部项消失，所以它们充分。如果一个端点可在曲面内任意移动，长度的边界项便不可能对所有向量都为零，因为 \(T\) 是非零切向量。这说明最短路问题必须先说清楚端点允许怎样动。

\subsection{测地坐标和径向局部极小性}
\label{g1:subsec:0902:04}

第一变分只给出了极小的必要条件；要证明短测地线确实最短，还需要一种能同时比较许多竞争曲线的坐标。直觉是从同一点向各个方向发出测地线，用初速度标记它们的终点。这一标记是否会重合，必须先在局部解决。

\begin{definition}
	\label{g1:def:surface-exponential}
	设 \((M,g)\) 为光滑度量曲面，\(p\) 为其内点。对 \(v\in\mathrm T_pM\)，以 \(\gamma_v\) 记满足 \(\gamma_v(0)=p\)、\(\dot\gamma_v(0)=v\) 的测地线。对使 \(\gamma_v\) 存在于 \([0,1]\) 的 \(v\in \mathrm{T}_pM\)，定义\term[zhishuyingshe]{指数映射}[exponential map]
	\[
		\exp_p(v)=\gamma_v(1).
	\]
	这个名称表示由初速度生成测地线的映射，并不是坐标逐项取通常的指数函数。\symidx{\(\exp_p\)：曲面的指数映射}
\end{definition}

\begin{proposition}[正常邻域]
	\label{g1:prop:normal-neighborhood}
	设 \(M\) 为光滑度量曲面，\(p\) 为内点。则 \(\exp_p\) 在 \(\mathrm{T}_pM\) 中零向量的某个开邻域上光滑，并满足 \(\mathrm{d}(\exp_p)_0=\operatorname{id}_{\mathrm{T}_pM}\)。存在 \(\rho>0\)，使它在 \(B_{2\rho}(0)\subset \mathrm{T}_pM\) 上为到开邻域的光滑微分同胚，其中球由 \(p\) 处的内积定义。
\end{proposition}
\begin{proof}
	零初速度给出常值测地线，在任意紧时间区间上都存在。由\cref{g1:lem:smooth-ode-flow}的有限时间结论，附近初速度的解也存在于包含 \([0,1]\) 的时间区间，并且光滑依赖初值，故 \(\exp_p\) 在零附近定义且光滑。

	初值唯一性给出缩放关系 \(\gamma_{tv}(s)=\gamma_v(ts)\)，只要两边有定义。因此对于固定 \(v\) 及足够小的 \(t\)，有 \(\exp_p(tv)=\gamma_v(t)\)，求 \(t=0\) 处的导数便得 \(\mathrm{d}(\exp_p)_0(v)=v\)。在 \(p\) 附近取一张坐标图后，\cref{g1:lem:surface-local-inverse}适用于这个二维光滑映射，给出零附近的光滑逆。最后在允许的邻域中选取足够小的 \(B_{2\rho}(0)\)。
\end{proof}

把 \(\mathrm{T}_pM\) 的一组正交单位基写作 \(e_1,e_2\)，用
\[
	(x^1,x^2)\longmapsto\exp_p(x^1e_1+x^2e_2)
\]
得到的坐标称为\term[cedizuobiao]{测地坐标}[geodesic normal coordinates]，其像称为\term[zhengchanglinyu]{正常邻域}[normal neighborhood]。这些坐标的优点是在中心点恢复 Euclid 度量的一阶形态。确实，\(\mathrm{d}(\exp_p)_0=\operatorname{id}\) 给出 \(g_{ij}(p)=\delta_{ij}\)；直线坐标 \(x(t)=tv\) 对应测地线，代入方程得 \(\Gamma^k_{ij}(p)v^iv^j=0\)。对所有 \(v\) 成立，加上 \(\Gamma^k_{ij}=\Gamma^k_{ji}\)，推出 \(\Gamma^k_{ij}(p)=0\)。度量相容的坐标式又给出 \(\partial_k g_{ij}(p)=0\)。这些消失只发生在中心点，不能据此认定整个邻域平坦。

\begin{lemma}[Gauss 引理]
	\label{g1:lem:gauss-lemma}
	设 \(p\) 是光滑度量曲面的内点，\(\exp_p\) 在零点的星形开邻域 \(B\subset \mathrm{T}_pM\) 上光滑且其中所有径向测地线均定义至时间 \(1\)。则对 \(v\in B\) 及 \(w\in \mathrm{T}_pM\)，有
	\[
		\left\langle\mathrm{d}(\exp_p)_v(v),\mathrm{d}(\exp_p)_v(w)\right\rangle_{\exp_p(v)}
		=\langle v,w\rangle_p.
	\]
	特别地，指数映射保持径向长度，并把与径向垂直的向量映为与径向测地线垂直的向量。
\end{lemma}
\begin{proof}
	\(v=0\) 时两端均为零。其余情况取小 \(s\)，令 \(F(t,s)=\exp_p(t(v+sw))\)、\(0\le t\le1\)。星形条件保证这些线段留在定义域内。记 \(A=\partial_tF\)、\(B_1=\partial_sF\)。每条 \(t\) 曲线为测地线，故 \(\mathrm{D}_tA=0\)，且速度平方为 \(|v+sw|_p^2\)。利用混合协变导数相等，得到
	\[
		\partial_t\langle A,B_1\rangle
		=\langle A,\mathrm{D}_sA\rangle
		=\frac12\partial_s|A|^2
		=\langle v+sw,w\rangle_p.
	\]
	在 \(t=0\) 处，\(F(0,s)=p\)，所以 \(B_1(0,s)=0\)。积分到 \(t=1\) 并取 \(s=0\)，左边正是所求的内积。令 \(w=v\) 给出径向长度保持，令 \(w\perp v\) 给出径向正交性。
\end{proof}

这一结论称为\term[gaussyinli]{Gauss 引理}[Gauss lemma]。它没有说指数映射保持所有长度。事实上，取 \(e(\theta)=e_1\cos\theta+e_2\sin\theta\)，在 \(r>0\) 处使用极坐标
\[
	(r,\theta)\longmapsto\exp_p(re(\theta)),
\]
Gauss 引理只保证径向速度长度为 \(1\)、径向与角向正交。度量因而具有形式
\begin{equation}
	\label{g1:eq:geodesic-polar-metric}
	I=\mathrm{d}r^2+a(r,\theta)^2\,\mathrm{d}\theta^2,\qquad a(r,\theta)>0.
\end{equation}
其中 \(a(r,\theta)=|\mathrm{d}(\exp_p)_{re(\theta)}(r e'(\theta))|\)。光滑性给出 \(a/r\to1\)，但一般 \(a\ne r\)。例如球面中心为 \(p\) 时，\cref{g1:exm:sphere-geodesics}给出 \(a=R\sin(r/R)\)，在 \(0<r<\pi R\) 内有定义，到了对跖点则所有径向方向相遇。极坐标在 \(r=0\) 的角度没有定义，光滑的测地直角坐标却没有这个缺陷。

% !TeX root = ../Book1.tex
\begin{figure}[htbp]
	\centering
	\begin{tikzpicture}[scale=0.95,font=\small]
		\begin{scope}[xshift=-3cm]
			\draw[gray,dashed] (0,0) circle (1.6);
			\draw[blue!70!black,thick,->] (0,0)--(1.05,0);
			\draw[red!70!black,thick,->] (0,0) .. controls (-0.1,1.05) and (1.4,1.2) .. (1.05,0);
			\fill (0,0) circle (1.6pt) node[below left] {\(p\)};
			\fill (1.05,0) circle (1.6pt) node[below right] {\(q\)};
			\node at (0,-2) {留在 \(V\) 中的曲线};
		\end{scope}
		\begin{scope}[xshift=3cm]
			\draw[gray,dashed] (0,0) circle (1.6);
			\draw[blue!70!black,thick,->] (0,0)--(1.05,0);
			\draw[red!70!black,thick,->] (0,0) .. controls (-0.4,1.4) and (0,2.3) .. (1.4,2) .. controls (2.7,1.7) and (2.1,-0.15) .. (1.05,0);
			\fill (0,0) circle (1.6pt) node[below left] {\(p\)};
			\fill (1.05,0) circle (1.6pt) node[below right] {\(q\)};
			\node at (0,-2) {先到达 \(r=\rho\) 的曲线};
		\end{scope}
	\end{tikzpicture}
	\caption[正常邻域中的最短路比较]{测地坐标中的两类竞争曲线。蓝线的长度为 \(r(q)\)；左图中角向运动会增加长度，右图的曲线仅走到虚线边界就至少需要长度 \(\rho\)。图中纸面上的一般距离不代表曲面距离。}
	\label{g1:fig:normal-neighborhood}
\end{figure}


\begin{theorem}[径向测地线的最短性]
	\label{g1:thm:radial-minimizing}
	设 \((M,g)\) 是光滑度量曲面，\(p\) 为内点，并取 \(\rho>0\)，使 \(\exp_p\) 在 \(B_{2\rho}(0)\) 上为微分同胚。令 \(V=\exp_p(B_\rho(0))\)。对 \(q=\exp_p(v)\in V\)，曲线 \(t\mapsto\exp_p(tv)\)、\(0\le t\le1\)，在 \(M\) 内所有连接 \(p,q\) 且参数化在 \([0,1]\) 上的分片 \(C^1\) 曲线中长度最小，长度为 \(|v|\)。当 \(q\ne p\) 时，达到该长度的正则分片光滑曲线具有同一径向轨迹，并沿同一方向行进；在共同区间 \([0,1]\) 上以恒速参数化后，最短参数曲线唯一。
\end{theorem}
\begin{proof}
	记 \(r(x)=|\exp_p^{-1}(x)|\)。先比较留在较大正常邻域中的曲线。Gauss 引理把非中心点处任意速度 \(W\) 正交分解为径向部分与角向部分，因而
	\[
		|W|^2=\bigl(\mathrm{d}r(W)\bigr)^2+|W_{\mathrm{ang}}|^2
		\ge\bigl(\mathrm{d}r(W)\bigr)^2.
	\]
	这也可以直接从\eqref{g1:eq:geodesic-polar-metric}读出。设竞争曲线 \(\eta:[a,b]\to V\) 从 \(p\) 到 \(q\ne p\)。为避开中心处的极坐标问题，对任意 \(0<\varepsilon<r(q)\)，取最后一个满足 \(r(\eta(t_\varepsilon))=\varepsilon\) 的时刻。连续性及终点条件保证它存在，且随后 \(r>\varepsilon\)。在这一段上逐段积分，得到
	\[
		L(\eta)\ge\int_{t_\varepsilon}^b|(r\circ\eta)'(t)|\,\mathrm{d}t
		\ge r(q)-\varepsilon.
	\]
	令 \(\varepsilon\downarrow0\)，便有 \(L(\eta)\ge r(q)\)。径向曲线的速度恒为 \(|v|\)，所以恰好达到这个界。

	现在允许竞争曲线离开 \(V\)。由于 \(\exp_p\) 在更大的球 \(B_{2\rho}\) 上是微分同胚，\(\overline V=\exp_p(\overline B_\rho)\) 为紧集，其边界就是 \(\exp_p(\partial B_\rho)\)。任何从 \(p\) 出发又离开 \(V\) 的连续曲线，必有一个首先到达该边界的时刻；此前的弧段仍处于正常坐标内。刚才的估计对这一弧段给出长度至少为 \(\rho\)，而 \(r(q)<\rho\)。因此离开 \(V\) 的曲线不可能达到径向曲线的长度。

	最后说明等号。先注意，位于 \(V\) 的任意弧段的长度至少为两端的径向坐标差的绝对值：不经过 \(p\) 时直接积分 \(|(r\circ\eta)'|\)；经过 \(p\) 时，取第一次及最后一次到达 \(p\) 的时刻，分别把两端弧段按上述从中心出发的估计处理，便得到更强的下界 \(r(\eta(a))+r(\eta(b))\)。因此，若总长度恰为 \(r(q)\)，在曲线不经过 \(p\) 的任意一小段上，径向增量的积分不能严格小于长度：否则把这段的严格差值与前后两段的径向估计相加，会使总长度严格大于 \(r(q)\)。于是该段必须同时满足 \(W_{\mathrm{ang}}=0\) 及 \((r\circ\eta)'\ge0\)。正则性使离开 \(p\) 后的径向速度为正，角方向保持常数；它也不可能返回 \(p\)，否则已经走过的正长度不能被后续消去。因此轨迹为从 \(p\) 沿同一方向到 \(q\) 的径向段。其恒速参数由端点、参数区间和长度唯一确定。
\end{proof}

最短性是以初始点为中心、在足够小范围内成立的结论。正常邻域中的任意两点未必能直接用这一中心点的径向线连接；要讨论它们之间的唯一最短线，须重新以其中一点作上述构造，并核对另一点落在允许范围内。球面大圆可以继续延伸，超过半圈后却不再最短。初值方程控制测地线的延伸，Gauss 引理则给出短段的长度比较；延伸范围与最短范围由此成为两个问题。

% !TeX root = ../../Book1.tex
% 纲要标题：### 9.3 局部 Gauss–Bonnet 公式
% 纲要位置：几何学计划纲要.md，第 479 行。
% 正文状态：初稿；尚未正式审核。

% 纲要内容开始
% 
% * 切向转角与联络积分
% * 分片光滑边界及外角
% * 圆盘区域上的公式与定向符号
% * 测地三角形的角亏
% 
% 纲要内容结束

\section{局部 Gauss–Bonnet 公式}
\label{g1:sec:0903}

平面上的正向简单闭曲线，其切向量净转过一周。曲面上的切向量却分属不同的切平面，必须先说明用什么方向作参照。联络给出比较方向的方法，测地曲率记录相对于平行移动的转向；把它沿边界积分，通常不再得到 \(2\pi\)。本节将说明，恰好缺少的量是区域内部的 Gauss 曲率积分。边界若有角点，还须计入转弯时发生的方向跳跃。

\subsection{切向转角与联络积分}
\label{g1:subsec:0903:01}

设曲面已经定向，\(J\) 表示切平面内正向旋转 \(\pi/2\)。取正向正交单位标架 \((e_1,e_2)\)，所以 \(Je_1=e_2\)、\(Je_2=-e_1\)。回忆\eqref{g1:eq:surface-frame-connection}中的约定：联络一形式 \(\omega\) 满足
\[
	\nabla_Ve_1=\omega(V)e_2,\qquad \nabla_Ve_2=-\omega(V)e_1.
\]
沿一段按弧长参数化的 \(C^2\) 正则曲线 \(\gamma\)，单位切向量 \(T=\gamma'\) 可写成
\[
	T(s)=\cos\theta(s)e_1+\sin\theta(s)e_2.
\]
这里标架在 \(\gamma(s)\) 处取值。\cref{g1:lem:circle-angle-lift}允许在参数区间上连续选择角度；因切向量是 \(C^1\) 的，局部光滑辐角使这个提升也是 \(C^1\) 的。角度不是曲面本身提供的绝对方向，而是切向量相对于所选标架的方向。

\begin{proposition}[切向转角公式]
	\label{g1:prop:tangent-connection-angle}
	设 \(\gamma\) 是有向光滑曲面上的 \(C^2\) 弧长参数正则曲线，\((e_1,e_2)\) 是其邻域中的光滑正向正交单位标架，\(\omega=\langle\nabla e_1,e_2\rangle\)。若 \(T=\cos\theta\,e_1+\sin\theta\,e_2\)，则
	\begin{equation}
		\label{g1:eq:tangent-connection-angle}
		k_g=\theta'+\omega(T),\qquad
		\int_\gamma k_g\,\mathrm{d}s=\theta(b)-\theta(a)+\int_\gamma\omega.
	\end{equation}
\end{proposition}
\begin{proof}
	按\cref{g1:prop:along-derivative-properties}的乘法法则沿曲线求协变导数，得到
	\begin{align*}
		\mathrm{D}_sT
		&=-\theta'\sin\theta\,e_1+\theta'\cos\theta\,e_2
		+\cos\theta\,\omega(T)e_2-\sin\theta\,\omega(T)e_1\\
		&=\bigl(\theta'+\omega(T)\bigr)JT.
	\end{align*}
	\cref{g1:def:geodesic-curvature}规定 \(k_g=\langle \mathrm{D}_sT,JT\rangle\)，故得第一式。对 \(s\) 积分，并用一形式的线积分定义 \(\int_\gamma\omega=\int_a^b\omega(T)\,\mathrm{d}s\)，便得第二式。
\end{proof}

这个等式把两种转动分开：\(\theta'\) 是切向量相对于标架的转动，\(\omega(T)\) 是标架自身相对于平行移动的转动。若沿曲线选择平行标架，第二项为零；若选择曲面上一组便于计算的坐标标架，第二项一般不能省略。尤其在曲线闭合时，沿它平行移动所得的标架未必首尾一致，因此不能同时把它当作周期标架，又任意令 \(\omega(T)=0\)。

\ProseParagraphBreak

换一组正向正交标架不会改变 \(k_g\)。例如局部令 \(\widetilde e_1=\cos\varphi\,e_1+\sin\varphi\,e_2\)、\(\widetilde e_2=J\widetilde e_1\)，则新角度是 \(\widetilde\theta=\theta-\varphi\)，新联络形式是 \(\widetilde\omega=\omega+\mathrm{d}\varphi\)。因而 \(\widetilde\theta'+\widetilde\omega(T)=\theta'+\omega(T)\)。计算中的角度和联络积分各自依赖参照，二者之和才是曲线的测地曲率积分。

\subsection{分片光滑边界及外角}
\label{g1:subsec:0903:02}

三角形的边可以是光滑的，但顶点处没有唯一切向量。若只把三条直线边的曲率积分相加，结果是零，整周转角便丢失了。因此要把沿边的连续转向和顶点处的瞬时转向一起记录。

\begin{definition}
	\label{g1:def:boundary-exterior-angle}
	设紧曲面区域 \(D\) 的边界由有限条 \(C^2\) 正则弧首尾相接而成，各弧在端点有非零单侧切向量，弧的内部互不相交。按区域位于左侧的方向走过边界，在角点 \(p_j\) 记到达和离开的单位切向量为 \(T_j^-\)、\(T_j^+\)。若区域在该点占据的切向扇形角为 \(\beta_j\in(0,2\pi)\)，称 \(\beta_j\) 为\term[neijiao]{内角}[interior angle]，并称
	\[
		\alpha_j\coloneqq\pi-\beta_j\in(-\pi,\pi)
	\]
	为该点的有向\term[waijiao]{外角}[exterior angle]。它是从 \(T_j^-\) 转到 \(T_j^+\) 的有向短转角；向左转取正，向右转取负。允许 \(\beta_j=\pi\)，此时只是人为分段点，外角为零。
\end{definition}

内角是由指向来路的 \(-T_j^-\) 与指向去路的 \(T_j^+\) 围成、朝向区域内部的角，不是两个前进切向量之间较小的夹角。\cref{g1:fig:boundary-exterior-angles}同时画出凸角与凹角：凹角的内角大于 \(\pi\)，外角便是负数。这里排除了内角为 \(0\) 或 \(2\pi\) 的尖点；在这种退化情形，两个前进切向量可能反向，不能直接使用唯一短转角。

\begin{figure}[htbp]
	\centering
	% !TeX root = ../Book1.tex
\begin{tikzpicture}[scale=1.05,>=stealth,font=\small]
	\begin{scope}
		\fill[blue!9] (0,0)--(60:1.18) arc[start angle=60,end angle=180,radius=1.18]--cycle;
		\draw[thick,->] (-1.75,0)--(-.12,0);
		\draw[thick,->] (0,0)--(60:1.65);
		\draw[dashed,gray,->] (0,0)--(1.2,0) node[right] {\(T^-\)};
		\draw[red!75!black,->] (.53,0) arc[start angle=0,end angle=60,radius=.53];
		\draw[blue!65!black] (60:.91) arc[start angle=60,end angle=180,radius=.91];
		\node at (1.15,.38) {\(\alpha>0\)};
		\node at (124:.66) {\(\beta<\pi\)};
		\node[above right] at (60:1.55) {\(T^+\)};
		\fill (0,0) circle (1.2pt);
		\node[below] at (0,-.16) {凸角：向左转};
	\end{scope}
	\begin{scope}[xshift=5cm]
		\fill[blue!9] (0,0)--(-60:1.18) arc[start angle=-60,end angle=180,radius=1.18]--cycle;
		\draw[thick,->] (-1.75,0)--(-.12,0);
		\draw[thick,->] (0,0)--(-60:1.65);
		\draw[dashed,gray,->] (0,0)--(1.2,0) node[right] {\(T^-\)};
		\draw[red!75!black,->] (.53,0) arc[start angle=0,end angle=-60,radius=.53];
		\draw[blue!65!black] (-60:.91) arc[start angle=-60,end angle=180,radius=.91];
		\node at (1.15,-.38) {\(\alpha<0\)};
		\node at (75:.63) {\(\beta>\pi\)};
		\node[below right] at (-60:1.55) {\(T^+\)};
		\fill (0,0) circle (1.2pt);
		\node[below] at (-.45,-1.55) {凹角：向右转};
	\end{scope}
\end{tikzpicture}

	\caption[边界外角的符号]{区域位于行进方向左侧。实线箭头表示先到达顶点、再沿下一边离开；虚线是到达方向的延长线。左图的内角小于 \(\pi\)，外角为正；右图的内角大于 \(\pi\)，外角为负。浅色扇形只示意顶点附近的内部方向。}
	\label{g1:fig:boundary-exterior-angles}
\end{figure}

将边界各段的切向角连续提升，并在每个角点接上转角 \(\alpha_j\)，便得到包含角点跳跃的总转角。这个总量应是多少，不能只从曲线闭合推出：闭合只使它成为 \(2\pi\) 的整数倍。下面说明，单图圆盘的正向简单边界使这个整数恰好为 \(1\)。

\begin{lemma}[分片边界的转角]
	\label{g1:lem:piecewise-turning}
	设有向光滑曲面上的紧区域 \(D\) 同胚于闭圆盘，并且包含在一张正向坐标图中。假设 \(\partial D\) 是正向分片 \(C^2\) 正则简单闭曲线，各角点内角均属于 \((0,2\pi)\)。将该图的坐标向量按顺序正交单位化，得到正向标架 \((e_1,e_2)\)。若 \(\Delta\theta_i\) 是第 \(i\) 条边的切向量相对于该标架的角增量，\(\alpha_j\) 为各外角，则
	\begin{equation}
		\label{g1:eq:piecewise-total-turning}
		\sum_i\Delta\theta_i+\sum_j\alpha_j=2\pi.
	\end{equation}
\end{lemma}
\begin{proof}
	先在坐标平面中证明带角点的转角定理。\cref{g1:thm:turning-tangent}已经证明 \(C^1\) 正则简单闭平面曲线的正向总转角为 \(2\pi\)；现在只需把角点的贡献补回来。在某个角点，两个前进单位切向量不反向，取它们的角平分方向为横轴，则两个向量在该横轴上的分量都为正。于是足够短的两侧曲线合在一起可以写成连续图像 \(y=f(x)\)，其左右单侧导数存在，且 \(f\) 在两侧为 \(C^2\)。取互不相交的小邻域处理全部角点，并使这些邻域避开其余曲线。

	在一个这样的邻域中，把 \(f\) 在角点附近换成光滑近似，并用光滑截断函数使它在邻域边缘与原函数完全相同。这个操作可以具体实施为：先在稍大的区间上延拓 \(f\)，用非负光滑积分核作卷积，再在角点附近取卷积结果、远离角点取原函数，在两者之间以截断函数相接。卷积在小尺度下一致趋于连续函数 \(f\)，所以改动后的图像仍留在所选邻域。它始终是关于 \(x\) 的图像，故不产生自交，也不会使前进切向量为零。所得闭曲线为 \(C^2\) 正则简单曲线，并沿用原曲线的正向。

	图像段的切向角可取为 \(\arctan f'(x)\)，值始终在 \((-\pi/2,\pi/2)\) 内；因此，无论中间怎样光滑化，这一小段的总转角恰为其终端与初端的切向角之差。让修改区间缩向原角点，两端角度趋向单侧切向角，差趋向该角点的有向短转角。未修改各段的角增量则趋向原边的角增量。将光滑曲线的转角定理取极限，得到坐标平面中“沿边角增量加外角等于 \(2\pi\)”。内角为 \(\pi\) 的分段点不改变这一论证。

	还需比较坐标平面的角和曲面度量测得的角。记坐标为 \((u,v)\)，度量系数为 \(E,F,G\)，\(D_g=EG-F^2>0\)。正交单位化给出
	\[
		\partial_u=\sqrt E\,e_1,\qquad
		\partial_v=\frac{F}{\sqrt E}\,e_1+\sqrt{\frac{D_g}{E}}\,e_2.
	\]
	因此，坐标速度 \(w=(\dot u,\dot v)^{\mathsf T}\) 在正交标架中的分量是
	\[
		Pw,\qquad
		P=\begin{pmatrix}\sqrt E&F/\sqrt E\\0&\sqrt{D_g/E}\end{pmatrix}.
	\]
	对 \(0\leq t\leq1\)，矩阵 \(P_t=(1-t)I+tP\) 都是对角元为正的上三角矩阵，故可逆且保向。沿每条边将 \(P_tw\) 归一化，在每个角点将左右极限用有向短圆弧接起来，就得到方向圆周上的连续闭路。这个闭路随 \(t\) 连续变化：非反向的两个端点经过可逆线性变换后仍不反向，故短圆弧不会突然跳到另一边。可给每个角点插入一段固定长度的辅助参数区间，使这里的闭路具有同一个参数圆周。

	\cref{g1:thm:circle-fundamental}的绕数同伦不变性说明，这个方向闭路在 \(t=0\) 与 \(t=1\) 时有相同总角增量。\(t=0\) 是刚证明的平面分片边界，总增量为 \(2\pi\)；\(t=1\) 正是曲面切向量相对于 \((e_1,e_2)\) 的方向及度量外角。这就得到\eqref{g1:eq:piecewise-total-turning}。
\end{proof}

平面角和曲面角一般不相等；上面保持的是绕一周的整数，而不是每个顶点的数值。这也解释了为何不能直接把坐标图画出的三角形内角，当成曲面上的内角代入公式。

\subsection{圆盘区域上的公式与定向符号}
\label{g1:subsec:0903:03}

现在已有两个可以相接的等式：切向转角公式把边界的测地曲率写成角增量与联络积分之和，曲率二形式则把联络的外微分写成 \(-K\,\mathrm{d}A\)。正号、负号都由这两个约定决定，不另作猜测。

\begin{theorem}[局部 Gauss–Bonnet 公式]
	\label{g1:thm:local-gauss-bonnet}
	设 \(M\) 是带光滑正定度量的有向光滑曲面，\(D\subset M\) 是包含在一张正向坐标图中的紧闭圆盘区域。假设其边界是有限条 \(C^2\) 正则弧组成的简单闭曲线，所有内角 \(\beta_j\) 均在 \((0,2\pi)\) 内。边界按 \(D\) 位于左侧定向，\(k_g=\langle \mathrm{D}_sT,JT\rangle\)，\(\alpha_j=\pi-\beta_j\)，则
	\begin{equation}
		\label{g1:eq:local-gauss-bonnet}
		\int_D K\,\mathrm{d}A
		+\int_{\partial D}k_g\,\mathrm{d}s
		+\sum_j\alpha_j=2\pi.
	\end{equation}
	边界积分表示各光滑边上积分之和；没有角点时，最后一项为空和。
\end{theorem}
\begin{proof}
	在给定坐标图上将坐标向量正交单位化，得到覆盖整个 \(D\) 及其邻域的正向标架，记其联络形式为 \(\omega\)。对每条边使用\eqref{g1:eq:tangent-connection-angle}，加上所有外角，再用\cref{g1:lem:piecewise-turning}，得到
	\[
		\int_{\partial D}k_g\,\mathrm{d}s+\sum_j\alpha_j
		=2\pi+\int_{\partial D}\omega.
	\]
	由\cref{g1:lem:green-formula}，当前分片 \(C^2\) 边界及正向参数满足二维积分公式的条件，故
	\[
		\int_{\partial D}\omega=\int_D\mathrm{d}\omega
		=-\int_D K\,\mathrm{d}A,
	\]
	最后一步是\cref{g1:prop:connection-curvature-form}的曲率形式公式。代回即得结论。
\end{proof}

这就是带边界和角点的\term[jubugaussbonnetgongshi]{局部 Gauss–Bonnet 公式}[local Gauss--Bonnet formula]。曲面可以嵌入三维欧氏空间，也可以只给定二维度量；证明仅用了度量的协变导数及其曲率。区域是拓扑圆盘，并不自动意味着它被当前某张坐标图覆盖。本定理明确保留单图假设；\secref{g1:sec:0904}将通过分解与粘合去掉它，并处理有孔的区域。

\ProseParagraphBreak

公式中的 \(\mathrm{d}A\) 在正向坐标下就是正的面积测度。若保持曲面定向，却把边界反向走，\(k_g\) 和有向转角都变号；以 \(-\partial D\) 表示这个反向边界，正确的式子是
\[
	-\int_D K\,\mathrm{d}A+\int_{-\partial D}k_g\,\mathrm{d}s+\sum_j(-\alpha_j)=-2\pi.
\]
这里仍按正面积测量 \(D\)，所以面积积分前也须带负号。若将曲面定向与边界方向一起反转，则 \(J\) 与 \(T\) 同时变号，\(\mathrm{D}_sT\) 保持不变，故 \(\langle \mathrm{D}_sT,JT\rangle\) 不变；区域内角与 \(K\) 也不变，原公式便仍成立。正向边界上的 \(JT\) 是指向区域内部的单位余法向，因而此时 \(k_g\) 也可理解为边界加速度的向内分量。

\begin{example}[平面圆盘与多边形]
	\label{g1:exm:planar-gauss-bonnet}
	在平面圆盘上，\(K=0\)，半径 \(r\) 的正向圆周满足 \(k_g=1/r\)、长度为 \(2\pi r\)，因此边界积分为 \(2\pi\)。对有 \(m\) 个顶点的平面简单多边形，各边是直线，\(k_g=0\)，所以公式给出
	\[
		\sum_{j=1}^m(\pi-\beta_j)=2\pi,
		\qquad \sum_{j=1}^m\beta_j=(m-2)\pi.
	\]
	这不要求多边形凸；凹顶点必须采用大于 \(\pi\) 的内部角，外角相应取负。
\end{example}

\subsection{测地三角形的角亏}
\label{g1:subsec:0903:04}

平面三角形的边没有沿平面的弯曲，因此全部转向都集中在顶点。在曲面上，测地线正好承担同样的角色：其单位切向量平行，故 \(\mathrm{D}_sT=0\)，从而 \(k_g=0\)。把三条测地线弧组成的简单边界代入局部公式，就能直接读出三角形的角和。

\begin{corollary}[测地三角形的角和]
	\label{g1:cor:geodesic-triangle-defect}
	设有向光滑曲面上一块单图闭圆盘区域 \(D\) 的边界由三条正则测地线弧组成，弧的内部互不相交，三个内角 \(\beta_1,\beta_2,\beta_3\) 均属于 \((0,2\pi)\)。则
	\begin{equation}
		\label{g1:eq:geodesic-triangle-defect}
		\beta_1+\beta_2+\beta_3-\pi=\int_D K\,\mathrm{d}A.
	\end{equation}
	特别地，若 \(K\equiv c\) 在 \(D\) 上为常数，则角和为 \(\pi+c\operatorname{Area}(D)\)。
\end{corollary}
\begin{proof}
	沿每条测地线取弧长参数后，测地线方程给出 \(\mathrm{D}_sT=0\)，所以三条边的测地曲率积分全为零。由\eqref{g1:eq:local-gauss-bonnet}，
	\[
		\int_D K\,\mathrm{d}A+3\pi-(\beta_1+\beta_2+\beta_3)=2\pi.
	\]
	移项即得结论。常曲率情形把常数提出面积积分即可。
\end{proof}

称 \(\delta=\pi-(\beta_1+\beta_2+\beta_3)\) 为三角形的\term[jiaokui]{角亏}[angular defect]；因此本约定下 \(\delta=-\int_DK\,\mathrm{d}A\)。正曲率给出角盈，负曲率给出正的角亏；若曲率变号，则由整个区域的积分决定，不能只看某个顶点的曲率。

\begin{example}[球面上的三个直角]
	\label{g1:exm:spherical-octant}
	在半径 \(R\) 的球面上，取三个点 \(A=(R,0,0)\)、\(B=(0,R,0)\)、\(C=(0,0,R)\)，用坐标平面中的四分之一大圆弧连接，围成第一卦限中的球面三角形。\cref{g1:exm:sphere-geodesics}已经给出大圆的测地性：用弧长参数化后，其空间加速度是 \(-\gamma/R^2\)，沿球面法向，故切向投影为零。每个顶点处两条边的切向量互相垂直，内角都是 \(\pi/2\)。

	这块闭三角形包含在开半球 \(x+y+z>0\) 中；沿法向 \((1,1,1)\) 的正交投影把该半球表示为球面的一个图像，故满足单图条件。\cref{g1:prop:surface-curvature-formulas}给出球面 \(K=1/R^2\)，于是
	\[
		\frac{\operatorname{Area}(D)}{R^2}
		=\frac{3\pi}{2}-\pi=\frac\pi2,
		\qquad \operatorname{Area}(D)=\frac{\pi R^2}{2}.
	\]
	这一结果也等于球面面积 \(4\pi R^2\) 的八分之一。三条边各自在曲面内都没有弯曲，三个直角仍能闭合，差别由区域内部的正曲率承担。
\end{example}

同样的代入还适用于 \(m\) 条测地线围成的单图圆盘：\(\sum_j\beta_j=(m-2)\pi+\int_DK\,\mathrm{d}A\)。但角和公式本身不保证给定三点之间能找到唯一的测地边，也不保证任意三条测地线围出圆盘。存在性、边的选择以及区域的拓扑条件仍须分别核对。

% !TeX root = ../../Book1.tex
% 纲要标题：### 9.4 整体公式的粘合证明
% 纲要位置：几何学计划纲要.md，第 486 行。
% 正文状态：初稿；尚未正式审核。

% 纲要内容开始
% 
% * 适应边界的曲面分解
% * 公共边项抵消及顶点贡献
% * Euler 示性数的出现
% * 闭曲面与带边界曲面的统一公式
% 
% 纲要内容结束

\section{整体公式的粘合证明}
\label{g1:sec:0904}

一个圆盘的局部公式右端是 \(2\pi\)，若把整张曲面切成 \(F\) 个圆盘，直接相加便得到 \(2\pi F\)。但切割也制造了原来没有的边和角：公共边的曲率会抵消，新增角的贡献却必须逐点数清。最后留下的整数并不是圆盘数，而是第六章的 \(V-E+F\)。这一过程还需要先解决一个正则性问题：第六章给出的是拓扑剖分，拓扑边没有可供积分的曲率，不能原样代入本章公式。

\subsection{适应边界的曲面分解}
\label{g1:subsec:0904:01}

下面允许 \(M\) 有分片 \(C^2\) 边界，角点有限，每个角点的内角属于 \((0,2\pi)\)。这里曲面和度量在各局部图中光滑；边界和角点表示在这样的光滑局部环境内截取一个区域，因而内角、边的切向量和协变加速度都有意义。内角为 \(\pi\) 的分段点也允许保留。

\begin{lemma}[坐标圆盘分解]\label{g1:lem:smooth-coordinate-disk-decomposition}
	设 \(M\) 是紧、Hausdorff、第二可数的光滑度量曲面，允许有分片 \(C^2\) 边界及有限个内角属于 \((0,2\pi)\) 的角点。则存在有限个闭圆盘 \(D_1,\ldots,D_F\)，满足：各 \(D_j\) 位于一张光滑坐标图中，边界由有限条 \(C^2\) 正则弧组成，各角的内角属于 \((0,2\pi)\)；圆盘的内部两两不交，合起来覆盖 \(M\)；两个圆盘的公共部分是若干条完整的边及顶点；\(\partial M\) 由分解的边和顶点组成，原有角点均为顶点。
\end{lemma}
\begin{proof}
	我们沿用\cref{g1:lem:finite-network-disk-cells}的构造思路：先造有限曲线网，再取顶点盘、边带，最后把平面区域中的孔连接起来。需要补充的是，所有新边都可以选成分片光滑的，并且每一块仍留在一张坐标图中。

	先在各点附近选带余量的小坐标片。在内点处用圆盘，在光滑边界点处用半圆盘；角点附近，在环境坐标中用小圆周截取两条边界弧围成的扇形区域。后一构造确实给出闭圆盘：两条正则弧从角点出发时，沿各弧的坐标半径在足够短的一段上严格增加，故每条小圆周恰与每侧相交一次；两侧的不同切向射线保证区域有正角。把最小各片的相对内部取成有限覆盖，再为每片保留稍大的同类片，其半径可以在一个开区间内变化而仍包含原来最小的片。这样的余量与\cref{g1:lem:nested-surface-charts}中的三层坐标片起同样作用。

	逐个选择这些较大片的半径，使新增的圆周弧与已经选定的边及 \(\partial M\) 有限次横交，并避开已有顶点。这个选择只需要如下一个实变量事实。若 \(f\in C^1([a,b])\)，则临界值集 \(f(\{f'=0\})\) 可以用总长度任意小的开区间覆盖。事实上，对任意 \(\varepsilon>0\)，在包含 \(\{f'=0\}\) 的开集中令 \(|f'|<\varepsilon\)；把该开集写成不交区间之并，平均值定理说明各像区间的长度不超过相应区间长度的 \(\varepsilon\) 倍，总和不超过 \(\varepsilon(b-a)\)。端点另用任意短的区间覆盖。因此有限个这类临界值集与有限个指定值，不能填满一个开区间。

	对每条已有光滑弧，令 \(f\) 为沿弧的坐标半径平方。只在所选的闭坐标环带上使用此函数；可以把径向函数在更大圆盘外取为常数，光滑延拓后再复合弧的参数，所以弧反复进出坐标图也不会妨碍上述论证。选择一个同时不是各弧临界值、也不是已有顶点半径平方的值。相交处沿弧的半径导数非零，正是两条曲线的切线不同；交点因而孤立。交点集又是紧集，故有限。避开已有顶点也排除了三条新旧边在同一点相交的情形。初始曲线网取为 \(\partial M\)，并把原有角点标出；最后给没有交点的闭曲线添上几个顶点，就得到一个有限的分片光滑嵌入图 \(G\)。

	每个 \(M\setminus G\) 的连通分支都留在某个已选坐标片内。确实，选分支中一点及包含它的坐标片内部；该片的相对边界已在 \(G\) 中，所以在不碰 \(G\) 的连通分支内，不能从片内走到片外。它的闭包也位于这张坐标图中，因为所选闭片与图的边缘之间仍有余量。不过补分支未必是盘，同一条边还可能在其边界行走中出现两次；下面的加宽步骤就是为了消除这种现象。

	在各顶点取两两不交的小坐标盘；边界顶点取相应半盘或角区域。半径取得足够小，每条关联边只以一个短弧进入顶点盘，并横交它的圆弧边界。把余下的边进一步有限分段，使每个闭段在一张坐标图中都是一条函数图像 \(v=f(u)\)；这是正则性与紧性给出的有限选择。在这样的图中，\(|v-f(u)|\le\delta\) 给出一条细带，位于 \(\partial M\) 上的边只取曲面内侧的半条带。可以令带的两端接到相邻顶点盘的边界弧上，并在接点附近调整侧边使其横交圆弧。有限条互不相交的截短边及其他不关联部分都是紧集，逐次缩小带宽即可使不同带的内部不交，且每个顶点盘、每条带都包含在一张坐标图中。顶点盘和带的并记为 \(N\)。它是 \(G\) 的闭邻域，包含 \(\partial M\) 的相对邻域，边界由有限条分片光滑 Jordan 曲线组成。所有接点都采用横交弧，故不会出现零角尖点。

	余下的每个闭区域 \(R\subset M\setminus\operatorname{int}_M N\) 位于原来一个补分支的坐标图中，其边界为有限条互不相交的分片光滑圈。区域只有有限多个：每个非空区域都接到 \(\partial N\) 的某条边界圈，而一条圈的外侧只邻接一个区域。\cref{g1:lem:finite-network-disk-cells}的平面部分说明，每个连通的 \(R\) 是一个外圈内去掉有限个内圈所围闭盘后的区域。若有内圈，在其一个非顶点处与外圈的非顶点处之间画一条内部位于 \(R^\circ\) 的简单分片光滑弧。具体地，先在两端取横交边界的小线段，再用 \(R^\circ\) 中的有限折线连接两段内端，删去折线回路，并在不碰边界的小圆盘内磨圆折线拐角。取这条弧的一条窄闭带，使其短边沿两条边界弧贴合。

	将窄带作为一个独立盘形块取出，剩下区域的两个边界圈合为一个，内圈数减少一；此处保留带的两个不同侧边，不把它们压回同一条割线。故余下区域始终具有真正的、互不相交的 Jordan 边界。有限次以后，每块余下区域只有一个边界圈，由\cref{g1:thm:jordan-schoenflies}是闭圆盘。它们和取出的窄带、顶点盘、原有边带一起给出有限闭盘分解，新增块都仍在所操作的原坐标图中。最后把所有接点统一作为顶点，边界上的分点也同时加入相邻两块。这样公共部分由完整边和顶点组成，原有边界和角点均保留下来。所有边由上述光滑弧构成，角可以选成严格介于 \(0\) 和 \(2\pi\) 之间，得到所需分解。
\end{proof}

这里没有要求分解的边是测地线。普通光滑边已经足够，因为它们的曲率积分会在粘合时抵消；若强求测地三角剖分，反而会额外引入尚未需要的存在性问题。闭盘也不必是三角形，下面按边和角的实际出现次数计数即可。

\subsection{公共边项抵消及顶点贡献}
\label{g1:subsec:0904:02}

先假设 \(M\) 可定向并选定定向，每块 \(D_j\) 的边界都按区域在左的方向行进。在一条内部公共边上，两块给出的行进方向相反。同一个曲线加速度在反向参数化后不变，而左侧的单位切向法向变号，因此两侧测地曲率互为相反数。两次积分使用相同的弧长测度，恰好抵消。

\ProseParagraphBreak

这个抵消其实不需要整张曲面可定向。对每块圆盘分别选一个局部定向，令其边界仍把本块放在左边。写 \(\nu_j\) 为沿边、垂直于边且指向 \(D_j\) 内侧的单位切向量，则正向行进时 \(\nu_j=JT\)，因而
\[
	k_{g,j}=\langle\nabla_TT,\nu_j\rangle.
\]
反向行进将 \(T\) 换成 \(-T\)，但 \(\nabla_{-T}(-T)=\nabla_TT\)；只要固定内侧，右端便不变。公共边两侧的 \(\nu_j\) 相反，所以无论各面选择的定向能否拼成整体定向，两项都抵消。沿原有边界，圆盘内侧与 \(M\) 内侧一致，留下的恰是 \(M\) 的边界曲率积分。于是
\begin{equation}\label{g1:eq:gb-edge-cancellation}
	\sum_{j=1}^F\int_{\partial D_j}k_{g,j}\,\mathrm{d}s
	=\int_{\partial M}k_g\,\mathrm{d}s.
\end{equation}
右端沿各条光滑边求和，不把角点另算成曲率函数的值。

\ProseParagraphBreak

再看角。用 \(V_{\mathrm i},V_{\mathrm b}\) 分别表示内部与边界顶点数，用 \(E_{\mathrm i},E_{\mathrm b}\) 表示内部与边界边数。每次一个圆盘在某顶点占据一个角，都算作一次角的出现，记其内角为 \(\beta_{j,v}\)。在内点，所有小角在切平面内围满一周，总和为 \(2\pi\)；在一个光滑边界点，它们铺满内侧半平面，总和为 \(\pi\)；在原有角点 \(v\)，总和等于 \(M\) 的内角 \(\beta_v\)，而不是 \(\pi\)。\cref{g1:fig:gb-vertex-angles}画的是顶点处带内积的切平面，各扇区就是这些角的切向模型。

% !TeX root = ../Book1.tex
\begin{figure}[htbp]
	\centering
	\begin{tikzpicture}[scale=0.88,font=\small]
		\begin{scope}
			\fill[GeometryBlueLight] (0,0) circle (1.15);
			\draw[GeometryBlue] (0,0) circle (1.15);
			\foreach \a in {0,55,135,210,285}{\draw[GeometryBlue] (0,0)--(\a:1.15);}
			\draw[GeometryBlue] (0.38,0) arc (0:55:0.38);
			\node at (27:0.7) {\(\beta_1\)};
			\fill (0,0) circle (1.3pt);
			\node[align=center] at (0,-1.65) {内点\\\(\sum\beta_j=2\pi\)};
		\end{scope}
		\begin{scope}[xshift=4cm]
			\fill[GeometryBlueLight] (-1.15,0)--(1.15,0) arc (0:180:1.15)--cycle;
			\draw[GeometryBlue] (-1.15,0)--(1.15,0) arc (0:180:1.15);
			\foreach \a in {45,110}{\draw[GeometryBlue] (0,0)--(\a:1.15);}
			\fill (0,0) circle (1.3pt);
			\node at (0,-0.45) {\(\partial M\)};
			\node[align=center] at (0,-1.65) {光滑边界点\\\(\sum\beta_j=\pi\)};
		\end{scope}
		\begin{scope}[xshift=8cm]
			\fill[GeometryBlueLight] (0,0)--(1.15,0) arc (0:125:1.15)--cycle;
			\draw[GeometryBlue] (0,0)--(1.15,0) arc (0:125:1.15)--cycle;
			\foreach \a in {40,83}{\draw[GeometryBlue] (0,0)--(\a:1.15);}
			\draw[GeometryBlue] (0.42,0) arc (0:125:0.42);
			\node at (63:0.72) {\(\beta_v\)};
			\fill (0,0) circle (1.3pt);
			\node[align=center] at (0.1,-1.65) {原有角点\\\(\sum\beta_j=\beta_v\)};
		\end{scope}
	\end{tikzpicture}
	\caption[粘合时的顶点角计数]{各面在同一顶点的切向扇区。内部顶点铺满一周，光滑边界顶点铺满半周，原有角点只铺满原内角。图中的圆弧仅用于标示扇区，不是曲面的分解边。}
	\label{g1:fig:gb-vertex-angles}
\end{figure}


为统一记号，给每个边界顶点置 \(\alpha_v=\pi-\beta_v\)，在新增的光滑分点处置 \(\alpha_v=0\)。则
\begin{equation}\label{g1:eq:gb-interior-angle-count}
	\sum_{j,v}\beta_{j,v}
	=2\pi V_{\mathrm i}+\pi V_{\mathrm b}
	-\sum_{v\in\partial M}\alpha_v.
\end{equation}
这里的角度是度量角，选择正交坐标后才与图中的 Euclid 角一致；角能相加是同一个切平面中的事实，不要求坐标图把曲面等距铺平。

\subsection{Euler 示性数的出现}
\label{g1:subsec:0904:03}

每个闭圆盘的边数等于角的出现次数。内部边出现两次，边界边出现一次，故总出现次数为 \(2E_{\mathrm i}+E_{\mathrm b}\)。把每个面角换成外角 \(\pi-\beta_{j,v}\)，再代入\eqref{g1:eq:gb-interior-angle-count}，得到
\begin{align*}
	\sum_{j,v}(\pi-\beta_{j,v})
	&=\pi(2E_{\mathrm i}+E_{\mathrm b})-2\pi V_{\mathrm i}-\pi V_{\mathrm b}
	+\sum_{v\in\partial M}\alpha_v\\
	&=2\pi(E_{\mathrm i}-V_{\mathrm i})+\sum_{v\in\partial M}\alpha_v.
\end{align*}
第二步用到 \(E_{\mathrm b}=V_{\mathrm b}\)：\cref{g1:prop:boundary-circles}说明紧曲面的边界是有限个圆周，每条圆周分成多少条边，就有多少个分点。这也解释了为何边界顶点不能与内部顶点一起按 \(2\pi\) 计角。

\ProseParagraphBreak

现在对每块 \(D_j\) 使用\cref{g1:thm:local-gauss-bonnet}。分解引理保证它是落在单张图中的闭圆盘，边界具有所要求的正则性和非零角，因此每块的右端都是 \(2\pi\)。对各块相加，面积积分可加、公共边按\eqref{g1:eq:gb-edge-cancellation}抵消，故
\[
	\int_MK\,\mathrm{d}A+\int_{\partial M}k_g\,\mathrm{d}s
	+\sum_{v\in\partial M}\alpha_v
	=2\pi(F-E_{\mathrm i}+V_{\mathrm i}).
\]
由于 \(E_{\mathrm b}=V_{\mathrm b}\)，右端括号等于
\[
	F-(E_{\mathrm i}+E_{\mathrm b})+(V_{\mathrm i}+V_{\mathrm b})=V-E+F=\chi(M).
\]
最后一步正是\cref{g1:prop:euler-triangulation-independent}：它不仅证明三角剖分的交错和不变，也证明任何有限闭盘分解的交错和都等于同一个 Euler 示性数。这里已经构造出真正的闭盘分解，因而满足该命题的条件。

\ProseParagraphBreak

可以把这段计数理解为两笔账。曲率积分记下各小片内部与边缘的几何；面片外角的总和则记下这些小片在顶点附近如何拼接。把人为切出的边和角消去后，几何账中留下总曲率，组合账中留下 \(\chi(M)\)。

\subsection{闭曲面与带边界曲面的统一公式}
\label{g1:subsec:0904:04}

\begin{theorem}[Gauss--Bonnet]\label{g1:thm:global-gauss-bonnet}
	设 \(M\) 为紧、Hausdorff、第二可数的光滑度量曲面，允许不连通、不可定向以及分片 \(C^2\) 边界。边界由有限条正则弧组成，角点的内角 \(\beta_v\) 属于 \((0,2\pi)\)。度量在边界及角点的局部光滑环境中光滑延拓。令 \(K\) 为 Gauss 曲率，\(\mathrm{d}A\) 为度量面积测度，沿每条边界弧取指向 \(M\) 内侧的单位切向法向 \(\nu\)，并定义
	\[
		k_g=\langle\nabla_TT,\nu\rangle,\qquad \alpha_v=\pi-\beta_v.
	\]
	这里 \(T\) 是沿该弧任一方向的单位切向量。则
	\begin{equation}\label{g1:eq:global-gauss-bonnet}
		\int_MK\,\mathrm{d}A+\int_{\partial M}k_g\,\mathrm{d}s+\sum_v\alpha_v
		=2\pi\chi(M).
	\end{equation}
	当 \(M\) 可定向时，边界按区域在左的方向行进，上述 \(k_g\) 就是 \(\langle\nabla_TT,JT\rangle\)。不连通时，\(\chi(M)\) 为各连通分支的 Euler 示性数之和。
\end{theorem}
\begin{proof}
	对每个连通分支使用\cref{g1:lem:smooth-coordinate-disk-decomposition}，按前两小节的边项抵消和顶点计数相加，即得\eqref{g1:eq:global-gauss-bonnet}。局部面积测度与 \(K\) 不因反转定向而改变；内侧测地曲率的抵消已经逐边证明，故不需要全局定向。最后，紧曲面只有有限个连通分支，这是\cref{g1:prop:boundary-circles}的另一项结论；将有限个公式相加即可。
\end{proof}

若边界光滑，角项为空；若 \(M\) 是闭曲面，边界积分也为空，公式成为
\begin{equation}\label{g1:eq:closed-gauss-bonnet}
	\int_MK\,\mathrm{d}A=2\pi\chi(M).
\end{equation}
另一方面，若 \(M\) 本身是一个圆盘，即使它不落在任何一张坐标图内，也有 \(\chi(M)=1\)，所以仍得到右端为 \(2\pi\) 的带边界公式。单图条件是局部证明所用的条件，经粘合后便不再是整体结论的限制。

\begin{example}\label{g1:exm:flat-mobius-gauss-bonnet}
	取带 Euclid 度量的矩形 \([0,L]\times[-a,a]\)，其中 \(L,a>0\)，按 \((0,y)\sim(L,-y)\) 识别两条竖边。因为识别映射在邻近条带上由刚性变换 \((x,y)\mapsto(x+L,-y)\) 给出，度量光滑粘合成一条 Möbius 带。这里的度量由矩形粘合定义，不需要把它作为 \(\mathbb R^3\) 中的曲面实现。

	局部图都是 Euclid 图，所以 \(K=0\)。上下两条水平边接成同一条光滑边界，各局部段都是直线，故 \(k_g=0\)，识别处也没有角点。\cref{g1:prop:surface-euler-formulas}给出 \(\chi(N_{1,1})=2-1-1=0\)，于是公式两边都是零。这个例子中不存在连续的全局定向，但“指向带内侧”的边界法向仍然存在，正好对应定理中采用的表述。
\end{example}

有边界时不能只看面积曲率。平面圆盘的 \(K\) 处处为零而 \(\chi=1\)，缺少的 \(2\pi\) 全由边界转动提供；平面环带的内圈则带来负的边界测地曲率，抵消外圈的贡献。角项承担同样的补偿作用，所以把光滑边界换成多边形边界后，公式的右端仍是同一个拓扑数。

% !TeX root = ../../Book1.tex
% 纲要标题：### 9.5 曲率积分的几何后果
% 纲要位置：几何学计划纲要.md，第 493 行。
% 正文状态：初稿；尚未正式审核。

% 纲要内容开始
% 
% * 亏格与曲率符号的限制
% * 总曲率与向量场零点的联系
% * 度量变化下的积分不变量
% 
% ---
% 
% 纲要内容结束

\section{曲率积分的几何后果}
\label{g1:sec:0905}

整体公式把一个可以逐点变化的函数，积分成了只能取离散值的拓扑量。因而它既能排除某些曲率分布，也能解释为什么改变度量时，总有一部分局部变化必须互相补偿。本节先看曲率符号，再把同一个积分理解为向量场绕零点的转动，最后计算保角改变度量时这种补偿怎样发生。

\subsection{亏格与曲率符号的限制}
\label{g1:subsec:0905:01}

本小节的曲面均紧、连通且无边界。第六章的\cref{g1:prop:surface-euler-formulas}给出：可定向亏格为 \(g\) 的曲面 \(\Sigma_g\) 满足 \(\chi=2-2g\)，不可定向标准曲面 \(N_k\) 满足 \(\chi=2-k\)，其中 \(k\ge1\)。因此\eqref{g1:eq:closed-gauss-bonnet}立即成为
\begin{equation}
	\label{g1:eq:total-curvature-genus}
	\int_{\Sigma_g}K\,\mathrm{d}A=4\pi(1-g),\qquad
	\int_{N_k}K\,\mathrm{d}A=2\pi(2-k).
\end{equation}
这里限制的是整张曲面的总积分。它并没有说一个点的曲率可以由亏格决定，也没有说总积分为零就意味着处处平坦。

\begin{proposition}
	\label{g1:prop:curvature-sign-topology}
	紧连通闭度量曲面满足下列限制。
	\begin{enumerate}
		\item 若 \(K\ge0\)，则 \(\chi\ge0\)；若还有一点 \(K>0\)，则 \(\chi>0\)。
		\item 若 \(K\le0\)，则 \(\chi\le0\)；若还有一点 \(K<0\)，则 \(\chi<0\)。
		\item 若 \(\chi=0\)，且 \(K\) 在全曲面上不变号，则 \(K\equiv0\)。
	\end{enumerate}
\end{proposition}
\begin{proof}
	积分保持非负性和非正性，故前两项的弱不等式来自 Gauss--Bonnet。若连续函数 \(K\) 在一点为正，则在某个小坐标圆盘上有 \(K\ge c>0\)。度量正定使这个圆盘的面积为正，所以在 \(K\ge0\) 的前提下，全曲率积分严格为正。负号情形同理。最后，若 \(K\) 不变号而又不恒为零，前两项就使积分严格非零，与 \(\chi=0\) 矛盾。
\end{proof}

结合\cref{g1:thm:surface-classification-unique}，这些限制可写成下表。表中列的是可能的拓扑类型，并不单凭这个必要条件断言某种指定曲率的度量一定存在。
\begin{center}
	\begin{tabular}{lll}
		\toprule
		曲率条件 & 可定向情形 & 不可定向情形\\
		\midrule
		\(K\ge0\)，且某处 \(K>0\) & 球面 \(\Sigma_0\) & 射影平面 \(N_1\)\\
		\(K\equiv0\) & 环面 \(\Sigma_1\) & Klein 瓶 \(N_2\)\\
		\(K\le0\)，且某处 \(K<0\) & \(\Sigma_g\)，\(g\ge2\) & \(N_k\)，\(k\ge3\)\\
		\bottomrule
	\end{tabular}
\end{center}

若 \(\Sigma_g\) 的曲率为常数 \(c<0\)，则还得到面积必须满足
\[
	\operatorname{Area}(\Sigma_g)=\frac{4\pi(g-1)}{-c}.
\]
所以固定曲率为 \(-1\) 后，闭可定向曲面的面积只能是 \(4\pi(g-1)\)。这个结论的前提仍是已经给定这样的度量；其构造属于后面的几何内容。

\begin{example}
	\label{g1:exm:two-torus-curvatures}
	平坦环面 \(\mathbb R^2/(a\mathbb Z\times b\mathbb Z)\) 的各坐标图通过平移相接，欧氏度量因而可以粘合。局部 Christoffel 系数全为零，故 \(K=0\)，总曲率当然为零。

	另一种熟悉的环面是三维空间中的旋转曲面
	\[
		x(u,v)=((R+r\cos v)\cos u,(R+r\cos v)\sin u,r\sin v),\qquad R>r>0.
	\]
	两个参数都以 \(2\pi\) 为周期。直接计算第一基本形式与形算子的行列式，得到
	\[
		\mathrm{d}A=r(R+r\cos v)\,\mathrm{d}u\,\mathrm{d}v,
		\qquad K=\frac{\cos v}{r(R+r\cos v)}.
	\]
	这个计算也可从\cref{g1:prop:revolution-two-forms}代入生成圆读出。外侧 \(\cos v>0\) 的曲率为正，内侧为负，且
	\[
		\int K\,\mathrm{d}A=\int_0^{2\pi}\int_0^{2\pi}\cos v\,\mathrm{d}v\,\mathrm{d}u=0.
	\]
	两种度量都符合环面的 Euler 示性数，但后者并不平坦。拓扑相同，也不意味着这两个度量局部等距。
\end{example}

\subsection{总曲率与向量场零点的联系}
\label{g1:subsec:0905:02}

一个处处非零的切向量场经单位化后，可以充当移动标架的第一根轴。若这样的轴只能在删掉某些点以后选取，就要问它绕着这些点转了几圈。第五章的圆周绕数正好能记录这个局部现象。

\begin{definition}
	\label{g1:def:surface-vector-index}
	设 \(V\) 是光滑切向量场，\(p\) 是其孤立零点。在 \(p\) 附近选一个定向坐标圆盘和正向光滑正交单位标架 \((e_1,e_2)\)，写 \(V=a e_1+b e_2\)。沿包围 \(p\)、内部没有其他零点的小正向坐标圆周 \(C\)，定义 \(V\) 在 \(p\) 的\term[xiangliangchangzhishu]{向量场指数}[index of a vector field]为映射
	\[
		C\longrightarrow S^1,\qquad q\longmapsto\frac{(a(q),b(q))}{\sqrt{a(q)^2+b(q)^2}}
	\]
	的绕数，记为 \(\operatorname{ind}_p(V)\)。等价地，若沿一周连续选择角度 \(\theta\)，则
	\[
		\operatorname{ind}_p(V)=\frac{\theta(\text{终点})-\theta(\text{起点})}{2\pi}.
	\]
	角度提升及所得整数的同伦不变性见\cref{g1:lem:circle-angle-lift,g1:thm:circle-fundamental}。\symidx{\(\operatorname{ind}_p(V)\)：切向量场在孤立零点的指数}
\end{definition}

这个整数与定义中的辅助选择无关。两条足够小的圆周在无零点的穿孔圆盘中可相互变形，单位化向量给出相应的圆周映射同伦。换正向正交标架时，新旧标架之间的旋转定义在整个圆盘上，其限制到边界的绕数为零；向量的角度减去这个旋转角，不改变一周的总增量。若改用另一张小坐标圆盘，可先在二者交中取更小的圆盘作比较。反转局部定向则同时反转边界行进方向和记录向量角度的方向，两个负号抵消。因此在不可定向曲面上，这个局部指数也有意义。

\begin{proposition}
	\label{g1:prop:nondegenerate-vector-index}
	在上述标架中置 \(F=(a,b)\)，并以零点为坐标原点。如果 \(A=\mathrm{d}F_0\) 可逆，则
	\[
		\operatorname{ind}_p(V)=\operatorname{sgn}\det A.
	\]
\end{proposition}
\begin{proof}
	由可逆性，存在 \(c>0\) 使 \(|Ax|\ge c|x|\)。Taylor 展开给出 \(F(x)=Ax+r(x)\)，其中 \(|r(x)|/|x|\to0\)。在足够小的圆周上，\(Ax+t r(x)\) 对 \(0\le t\le1\) 始终非零，所以 \(F/|F|\) 与 \(Ax/|Ax|\) 同伦。

	对 \(A\) 的列作 Gram--Schmidt 正交化，可写 \(A=QR\)，其中 \(R\) 上三角且对角元为正，\(Q\) 正交。矩阵 \((1-t)R+tI\) 始终可逆，故只需算 \(Q\) 的绕数。若 \(\det Q=1\)，它是旋转，绕数为 \(1\)；若 \(\det Q=-1\)，它是旋转与反射的复合，绕数为 \(-1\)。这正是 \(\operatorname{sgn}\det A\)。
\end{proof}

例如平面场 \((x,y)\) 与 \((-x,-y)\) 的指数都为 \(1\)：向外或向内改变的是箭头方向，不改变绕原点一周时箭头也转一周的事实。鞍点场 \((x,-y)\) 的指数是 \(-1\)。这一符号差别比“是不是零点”本身提供了更多信息。

\begin{theorem}[Poincar\'e--Hopf]
	\label{g1:thm:surface-poincare-hopf}
	设 \(M\) 为紧光滑闭曲面，\(V\) 为只有孤立零点的光滑切向量场。则零点只有有限个，且
	\begin{equation}
		\label{g1:eq:surface-index-sum}
		\sum_{V(p)=0}\operatorname{ind}_p(V)=\chi(M).
	\end{equation}
	允许 \(V\) 没有零点，此时左端为空和 \(0\)。
\end{theorem}

这是 Poincar\'e--Hopf 定理的二维情形。\footnote{Poincar\'e，庞加莱，Jules Henri Poincar\'e（1854--1912），法国数学家，对拓扑学、微分方程与动力系统作出奠基性贡献。}\footnote{Hopf，霍普夫，Heinz Hopf（1894--1971），德国出生的瑞士数学家，研究拓扑学、向量场及同伦不变量。}我们用刚证明的曲率公式给出证明，不预先调用高维版本。
\begin{proof}
	零点集闭且位于紧曲面中。每个零点都能被一个只含该零点的邻域隔开；若零点无限，紧性会给出一个聚点，连续性又使该聚点成为非孤立零点，矛盾。只需逐连通分支证明。

	先设 \(M\) 可定向。删去每个零点附近互不相交的坐标小开圆盘 \(D_p(\varepsilon)\)，所得曲面记为 \(M_\varepsilon\)。在它上面，
	\[
		E_1=\frac{V}{|V|},\qquad E_2=JE_1
	\]
	是一组全局正交单位标架，其联络形式 \(\Omega\) 满足 \(\mathrm{d}\Omega=-K\,\mathrm{d}A\)。Green 公式在这里也能逐块相加：用\cref{g1:lem:smooth-coordinate-disk-decomposition}分解 \(M_\varepsilon\)，对每块应用\cref{g1:lem:green-formula}，公共边上一形式的积分因方向相反而抵消。因此
	\[
		\int_{M_\varepsilon}K\,\mathrm{d}A=-\int_{\partial M_\varepsilon}\Omega.
	\]
	在 \(p\) 附近取一组能延到圆盘中心的正向光滑标架，联络形式记为 \(\omega_p\)。沿小圆周写 \(E_1=\cos\theta\,e_1+\sin\theta\,e_2\)。\cref{g1:prop:connection-frame-change}给出 \(\Omega=\omega_p+\mathrm{d}\theta\)。\(\theta\) 可以沿一周提升，不要求在穿孔圆盘上成为单值函数。按小圆盘的正向积分，便有
	\[
		\int_{\partial D_p(\varepsilon)}\Omega
		=\int_{\partial D_p(\varepsilon)}\omega_p+2\pi\operatorname{ind}_p(V).
	\]
	洞的边界在 \(M_\varepsilon\) 中取相反方向，故
	\[
		\int_{M_\varepsilon}K\,\mathrm{d}A
		=\sum_p\left(\int_{\partial D_p(\varepsilon)}\omega_p+2\pi\operatorname{ind}_p(V)\right).
	\]
	\(\omega_p\) 在整个小圆盘上光滑，其系数有界，圆周长度趋于零，因而前一项趋于零。删掉的圆盘面积也趋于零，而 \(K\) 有界。令 \(\varepsilon\to0\)，得到 \(\int_M K\,\mathrm{d}A=2\pi\sum_p\operatorname{ind}_p(V)\)，再用闭曲面的 Gauss--Bonnet 公式即得结论。如果没有零点，就直接在整个闭曲面上对 \(\Omega\) 积分，得到总曲率为零。

	若 \(M\) 不可定向，取\cref{g1:thm:orientation-cover}中的定向二重覆盖 \(\pi:\widetilde M\to M\)。以局部坐标的提升作为光滑坐标，并赋予拉回度量，\(\pi\) 就是光滑局部等距。场 \(V\) 唯一提升为 \(\widetilde V\)，每个零点提升成两个零点。局部微分同胚只是在同一指数定义中换坐标；即使它反转定向，前面解释的两个负号仍然抵消，所以两个指数都等于原指数。可定向情形给出 \(2\sum_p\operatorname{ind}_p(V)=\chi(\widetilde M)\)，再由\cref{g1:prop:surface-finite-cover-euler}知 \(\chi(\widetilde M)=2\chi(M)\)，除以 \(2\) 即得所求。
\end{proof}

\begin{example}
	\label{g1:exm:sphere-vector-indices}
	单位球面上高度函数 \(h(p)=p_3\) 的梯度是
	\[
		V(p)=e_3-p_3p.
	\]
	因为对每个切向量 \(w\) 都有 \(\langle V,w\rangle=w_3=\mathrm{d}h(w)\)，这个表达式确实给出曲面梯度。它仅在南、北极为零。北极处沿切方向求导得到 \(-I\)，南极处得到 \(I\)；二者在二维中行列式都为正，所以两个指数均为 \(1\)，总和为球面的 \(\chi=2\)。特别地，球面不可能有处处非零的光滑切向量场。平坦环面上的常向量场则能穿过平移识别而下降为处处非零的场，与 \(\chi=0\) 相符。
\end{example}

指数并不是给每个零点机械地计数一次。退化零点可以有更大的正指数、负指数或零指数；定理约束的是带指数的总和。它的证明还指出了联系所在：删去零点后标架存在，但标架绕这些洞的转角，恰好补回了整张曲面的曲率积分。

\subsection{度量变化下的积分不变量}
\label{g1:subsec:0905:03}

固定一张光滑闭曲面 \(M\)，任取两个光滑度量 \(g_0,g_1\)。虽然曲率和面积测度都依赖度量，Gauss--Bonnet 却给出
\[
	\int_MK_{g_0}\,\mathrm{d}A_{g_0}=\int_MK_{g_1}\,\mathrm{d}A_{g_1}=2\pi\chi(M).
\]
这里不需要先证明两种度量之间有什么特殊几何关系。事实上 \(g_t=(1-t)g_0+tg_1\) 对 \(0\le t\le1\) 始终正定，提供了一个直接的光滑变化；沿整个变化，总积分都固定。有边界时，固定的则是曲率积分、边界测地曲率积分与外角和三者之和。

为看清局部补偿，考虑保角改变度量 \(\widetilde g=e^{2f}g\)，其中 \(f\) 是光滑实函数。第八章在平面保角坐标中算过曲率；下面给出相对于任意曲面度量的公式。仍采用第八章约定 \(\Delta_g=-\operatorname{div}_g\operatorname{grad}_g\)。

\begin{proposition}
	\label{g1:prop:conformal-gauss-bonnet-change}
	在定向光滑度量曲面上，若 \(\widetilde g=e^{2f}g\)，则
	\begin{align}
		\mathrm{d}\widetilde A&=e^{2f}\mathrm{d}A,\label{g1:eq:conformal-area-change}\\
		\widetilde K\,\mathrm{d}\widetilde A&=(K+\Delta_g f)\,\mathrm{d}A.\label{g1:eq:conformal-curvature-change}
	\end{align}
	对区域的一段光滑边界，令 \(T\) 为原度量的正向单位切向量，\(\nu=JT\) 为内侧单位切向法向，\(\partial_\nu f=\mathrm{d}f(\nu)\)。则
	\begin{equation}
		\label{g1:eq:conformal-boundary-change}
		\widetilde k_g\,\mathrm{d}\widetilde s=(k_g-\partial_\nu f)\,\mathrm{d}s.
	\end{equation}
	保角变化保持每个角点的内角及外角。
\end{proposition}
\begin{proof}
	二维度量矩阵乘以 \(e^{2f}\) 后，其行列式乘以 \(e^{4f}\)，开方即得面积公式。为求曲率，先比较协变导数。定义
	\[
		\widetilde\nabla_VW=\nabla_VW+V(f)W+W(f)V-g(V,W)\operatorname{grad}_g f.
	\]
	右端新增项关于 \(V,W\) 对称，故仍无挠。对 \(\widetilde g(W,Z)=e^{2f}g(W,Z)\) 求导并展开，新增项在两个内积中的贡献为 \(2V(f)e^{2f}g(W,Z)\)，其余交叉项互相抵消，正好给出度量相容性。因此由\cref{g1:prop:surface-connection-properties}的唯一性，这就是新度量的协变导数。

	取原度量的正向正交单位标架 \((e_1,e_2)\)，以 \((\vartheta^1,\vartheta^2)\) 表示其对偶一形式。新标架为 \(\widetilde e_i=e^{-f}e_i\)。把它代入刚得到的公式，可得
	\[
		\widetilde\omega(V)=\omega(V)+e_1(f)g(V,e_2)-e_2(f)g(V,e_1).
	\]
	记 \(*\vartheta^1=\vartheta^2\)、\(*\vartheta^2=-\vartheta^1\)，线性延拓到一形式，则上述等式就是 \(\widetilde\omega=\omega+*\mathrm{d}f\)。这个星号这里只表示一形式在定向正交标架中的直角旋转；换正向正交标架不改变它的定义。

	还需要算出这个一形式的外微分。置 \(Z=\operatorname{grad}_g f=Z^u\partial_u+Z^v\partial_v\)，在正向坐标中写 \(h=\sqrt{\det(g_{ij})}\)。由定义
	\[
		*\mathrm{d}f=h(Z^u\mathrm{d}v-Z^v\mathrm{d}u),\qquad
		\mathrm{d}(*\mathrm{d}f)=\bigl(\partial_u(hZ^u)+\partial_v(hZ^v)\bigr)\mathrm{d}u\wedge\mathrm{d}v.
	\]
	第八章的散度公式表明，右端为 \((\operatorname{div}_g Z)\mathrm{d}A=-(\Delta_g f)\mathrm{d}A\)。因此
	\[
		-\widetilde K\,\mathrm{d}\widetilde A
		=\mathrm{d}\widetilde\omega
		=\mathrm{d}\omega+\mathrm{d}(*\mathrm{d}f)
		=-(K+\Delta_g f)\mathrm{d}A,
	\]
	即得曲率公式。

	沿边界，新旧标架中切向量的角度 \(\theta\) 相同，且 \(\mathrm{d}\widetilde s=e^f\mathrm{d}s\)。用\cref{g1:prop:tangent-connection-angle}将曲率乘弧长写成 \(\mathrm{d}\theta+\omega\)，得到
	\[
		\widetilde k_g\,\mathrm{d}\widetilde s
		=k_g\,\mathrm{d}s+(*\mathrm{d}f)(T)\,\mathrm{d}s
		=(k_g-\partial_\nu f)\,\mathrm{d}s.
	\]
	最后，度量乘以同一个正数不改变两条射线夹角的余弦，故角点的角也保持不变。
\end{proof}

曲率面积形式新增的 \(\Delta_g f\) 与边界新增项恰好抵消。具体地，由刚才的局部计算和逐块 Green 公式，
\[
	\int_M\Delta_g f\,\mathrm{d}A
	=-\int_{\partial M}*\mathrm{d}f
	=\int_{\partial M}\partial_\nu f\,\mathrm{d}s.
\]
这里使用的是内侧法向；若换成外侧，最后一项前面要加负号。把此式与\eqref{g1:eq:conformal-boundary-change}相加，总变化为零。闭曲面没有边界，便直接得到 \(\int_M\Delta_g f\,\mathrm{d}A=0\)。不可定向时，这些曲率和内侧法向公式在局部反转定向后仍不变，可以提升到定向二重覆盖再除以 \(2\)，所以总积分的结论照样成立。

\begin{example}
	\label{g1:exm:conformal-flat-torus}
	在 \(u,v\) 都以 \(2\pi\) 为周期的平坦环面上，取 \(f(u,v)=\varepsilon\cos u\)，其中 \(\varepsilon>0\)。新度量为 \(e^{2\varepsilon\cos u}(\mathrm{d}u^2+\mathrm{d}v^2)\)。因 \(K=0\)、\(\Delta_g f=-f_{uu}-f_{vv}=\varepsilon\cos u\)，有
	\[
		\widetilde K=\varepsilon e^{-2\varepsilon\cos u}\cos u,\qquad
		\widetilde K\,\mathrm{d}\widetilde A=\varepsilon\cos u\,\mathrm{d}u\,\mathrm{d}v.
	\]
	局部曲率已经出现正负变化，总积分却仍为零。这里曲率的补偿可以直接由周期积分看见；一般度量变化中，保证这种补偿的就是整体 Gauss--Bonnet 公式。
\end{example}

本章由切向投影出发，经联络形式、边界转角和有限分解，得到总曲率的拓扑不变性。向量场指数提供了同一整数的另一种计算方式。这些联系都已在二维中完成，不需要预先知道一般流形的特征类理论；后面的推广应当解释这里已经看清的现象，而不是替代这些具体证明。


\begin{chaptersummary}
	协变导数比较沿曲线移动的切向量；平行移动由线性常微分方程给出，并保持内积。在正向正交标架中，平行向量的角度满足 \(\theta'=-\omega(\dot\gamma)\)，而联络形式满足 \(\mathrm{d}\omega=-K\,\mathrm{d}A\)。因此绕一个坐标圆盘的平行移动，其旋转角等于内部曲率积分模 \(2\pi\) 的值。

	\ProseParagraphBreak

	测地线的速度沿自身平行。初值方程给出局部唯一解和光滑依赖，第一变分说明固定端点下的驻值条件；能量与长度之间还需要核对参数和恒速条件。指数映射与 Gauss 引理把径向和角向分开，在足够小的正常邻域内证明径向测地线最短。能继续延伸并不意味着能一直保持最短。

	\ProseParagraphBreak

	正向边界满足 \(k_g=\theta'+\omega(T)\)，角点另贡献 \(\alpha=\pi-\beta\)。圆盘的总切向转角为 \(2\pi\)，与 Green 公式相合便得到局部 Gauss--Bonnet。对有限坐标圆盘分解相加，内部边的测地曲率抵消，内部与边界顶点分别按完整角和内侧角计数，最终留下 \(2\pi(V-E+F)\)。内侧法向的表述使整体公式也适用于不可定向曲面。

	\ProseParagraphBreak

	闭曲面的总曲率等于 \(2\pi\chi\)，从而限制曲率符号与亏格。沿向量场的单位方向选标架，则零点附近的转动给出指数，总指数同样等于 \(\chi\)。改变度量可以改变每一点的曲率；在保角变化中，新增的曲率面积项为 \(\Delta_g f\,\mathrm{d}A\)，其积分由边界项补偿。总公式保持不变，局部几何仍可以发生显著变化。
\end{chaptersummary}

\BookExercises

\begin{exercise}\label{g1:ex:warped-annulus-holonomy}
	在环带 \(I\times(\mathbb R/2\pi\mathbb Z)\) 上给定度量 \(\mathrm{d}u^2+f(u)^2\mathrm{d}v^2\)，其中 \(I\) 为开区间，\(f>0\) 光滑，定向为 \(\mathrm{d}u\wedge\mathrm{d}v\)。
	\begin{enumerate}
		\item 对标架 \(e_1=\partial_u\)、\(e_2=f^{-1}\partial_v\)，求联络形式、Gauss 曲率及沿 \(u=u_0\) 增加 \(v\) 一周的平行移动旋转角。
		\item 取 \(I=(0,\infty)\)、\(f(u)=cu\)、\(c>0\)，说明曲率处处为零仍可能有非恒等的环路平行移动，并指出这为何不与局部曲率积分公式矛盾。
	\end{enumerate}
\end{exercise}
\begin{solution}
	1. 非零 Christoffel 系数为 \(\Gamma^u_{vv}=-ff'\)、\(\Gamma^v_{uv}=\Gamma^v_{vu}=f'/f\)。所以 \(\nabla_{\partial_u}e_1=0\)、\(\nabla_{\partial_v}e_1=f'e_2\)，给出 \(\omega=f'\mathrm{d}v\)。又有 \(\mathrm{d}A=f\,\mathrm{d}u\wedge\mathrm{d}v\)，故
	\[
		\mathrm{d}\omega=f''\mathrm{d}u\wedge\mathrm{d}v=-K\,\mathrm{d}A,
		\qquad K=-\frac{f''}{f}.
	\]
	这个标架沿 \(v\) 是周期的，平行向量相对它的角增量为 \(-\int\omega=-2\pi f'(u_0)\)。因此返回后的旋转角为该数模 \(2\pi\) 的值。

	2. 此时 \(K=0\)，旋转角却是 \(-2\pi c\pmod{2\pi}\)，当 \(c\notin\mathbb Z\) 时并非恒等。纬圈在环带内不围成圆盘，不能直接使用\cref{g1:prop:surface-local-holonomy}。把 \(u=0\) 添进去也不能随意补救：若 \(c\ne1\)，小圆周周长与径向距离之比是 \(2\pi c\)，不是光滑平面点应有的极限 \(2\pi\)，新添点具有锥形奇性。局部平坦性只保证足够小的可缩环路平行移动为恒等，不消除这种绕孔的整体效应。
\end{solution}

\begin{exercise}\label{g1:ex:sphere-octant-transport}
	在外向定向的单位球面上，依次沿短大圆弧从 \(N=(0,0,1)\) 到 \(A=(1,0,0)\)，再到 \(B=(0,1,0)\)，最后回到 \(N\)。把初始切向量 \((1,0,0)\in \mathrm{T}_NS^2\) 平行移动一周，直接求终点向量，并与所围球面八分之一的曲率积分比较。
\end{exercise}
\begin{solution}
	第一段可写 \(\gamma_1(t)=(\sin t,0,\cos t)\)，\(0\le t\le\pi/2\)。其单位切向量 \((\cos t,0,-\sin t)\) 的导数为 \(-\gamma_1(t)\)，纯属法向，所以它沿大圆平行。初始向量正是该切向量，到 \(A\) 时成为 \((0,0,-1)\)。

	第二段为赤道弧 \((\cos t,\sin t,0)\)。常向量 \((0,0,-1)\) 处处与球面法向垂直，环境导数为零，所以它平行地到达 \(B\)。第三段取 \(\gamma_3(t)=(0,\cos t,\sin t)\)，其单位切向量在 \(B\) 为 \((0,0,1)\)。待移动向量是该切向量的负值，故沿弧保持为 \((0,\sin t,-\cos t)\)，到 \(N\) 时成为 \((0,1,0)\)。北极的正向直角旋转把 \((1,0,0)\) 送到 \((0,1,0)\)，所以净旋转角为 \(\pi/2\)。

	该路径把第一卦限球面区域置于左侧；球面对坐标反射对称，八个卦限面积相同，总面积为 \(4\pi\)，故此块面积为 \(\pi/2\)。单位球面 \(K=1\)，曲率积分也为 \(\pi/2\)，与直接平行移动计算一致。该闭盘可放入一张避开其他球面点的坐标图，因此局部环路平行移动公式的条件满足；角点只分段改变行进方向，并不额外转动被移动的向量。
\end{solution}

\begin{exercise}\label{g1:ex:flat-rectangle-frame}
	设定向光滑度量曲面的一个坐标域为开矩形 \(U=(a,b)\times(c,d)\)，且 \(K\equiv0\)。不用一般的 Poincar\'e 引理，证明 \(U\) 上存在光滑正交单位标架，使两根标架向量沿任意曲线均平行。再证明任意两组这样的正向标架相差一个常旋转。
\end{exercise}
\begin{solution}
	把坐标基正交单位化，设所得联络形式为 \(\omega=A\mathrm{d}u+B\mathrm{d}v\)。曲率为零给出 \(B_u=A_v\)。固定 \((u_0,v_0)\in U\)，令
	\[
		h(u,v)=\int_{u_0}^uA(s,v_0)\,\mathrm{d}s+\int_{v_0}^vB(u,t)\,\mathrm{d}t.
	\]
	则 \(h_v=B\)，且
	\[
		h_u=A(u,v_0)+\int_{v_0}^v B_u(u,t)\,\mathrm{d}t
		=A(u,v_0)+\int_{v_0}^v A_v(u,t)\,\mathrm{d}t=A(u,v).
	\]
	所以 \(\mathrm{d}h=\omega\)。把原标架旋转角 \(-h\)，\cref{g1:prop:connection-frame-change}使新联络形式成为 \(\omega-\mathrm{d}h=0\)。两根新标架向量的协变导数因而在每个方向上都为零，沿任何曲线都平行。

	若 \((e_1,e_2)\) 与 \((f_1,f_2)\) 都有此性质，则对每个切向量 \(X\)，度量相容性给出 \(X\langle e_i,f_j\rangle=0\)。矩形连通，故四个内积都是常数。换基矩阵因此为常正交矩阵，两组标架均正向又使其行列式为 \(1\)，所以是常旋转。
\end{solution}

\begin{exercise}
	\label{g1:ex:cylinder-geodesics}
	在半径 \(R>0\) 的圆柱面上使用局部参数
	\[
		X(u,z)=(R\cos(u/R),R\sin(u/R),z).
	\]
	求出全部仿射参数测地线，判定哪些为闭曲线，并说明为什么闭合的一整圈不能是同端点之间的最短路。
\end{exercise}
\begin{solution}
	有 \(\langle X_u,X_u\rangle=\langle X_z,X_z\rangle=1\)、\(\langle X_u,X_z\rangle=0\)，所以度量为 \(\mathrm{d}u^2+\mathrm{d}z^2\)，全部 Christoffel 系数为零。测地线方程成为 \(\ddot u=\ddot z=0\)，因此
	\[
		\gamma(t)=\left(R\cos\frac{at+b}{R},R\sin\frac{at+b}{R},ct+d\right).
	\]
	这些公式在全部实数时间上有定义，并与每个局部初值相符，由初值唯一性给出全部测地线。\(a=0,c\ne0\) 时为竖直母线；\(a\ne0,c=0\) 时为水平圆；\(a,c\) 都非零时为圆柱螺线；\(a=c=0\) 时为常值曲线。非常值曲线闭合当且仅当 \(c=0\)，此时可取周期 \(2\pi R/|a|\)。走完一整圈的长度为 \(2\pi R>0\)，同端点之间的常值曲线长度为零，所以整圈不是最短路；这与其每个足够短的弧段为最短路并不冲突。
\end{solution}

\begin{exercise}
	\label{g1:ex:latitude-geodesic-curvature}
	半径 \(R>0\) 的球面取外向定向。对固定 \(0<\theta<\pi\)，按 \(\phi\) 增大的方向走纬线
	\[
		\gamma_\theta(\phi)=R(\sin\theta\cos\phi,\sin\theta\sin\phi,\cos\theta),\qquad 0\le\phi\le2\pi.
	\]
	计算其测地曲率，并用第一变分重新得到 \(\frac{\mathrm{d}}{\mathrm{d}\theta}L(\gamma_\theta)\)。这里闭曲线变分两端相接，端点向量不为零，需说明边界项为何仍相消。
\end{exercise}
\begin{solution}
	单位切向量和左侧单位向量分别为
	\begin{align*}
		T&=(-\sin\phi,\cos\phi,0),\\
		JT=n\times T&=(-\cos\theta\cos\phi,-\cos\theta\sin\phi,\sin\theta).
	\end{align*}
	由于 \(\mathrm{d}\ell=R\sin\theta\,\mathrm{d}\phi\)，空间中有
	\[
		\frac{\mathrm{d}T}{\mathrm{d}\ell}=\frac{(-\cos\phi,-\sin\phi,0)}{R\sin\theta}.
	\]
	与 \(JT\) 内积时法向投影不影响结果，故 \(k_g=\cot\theta/R\)。特别地，赤道 \(\theta=\pi/2\) 才是纬线族中的测地线。

	让 \(\theta\) 增大所给出的变分向量为 \(V=\partial_\theta\gamma_\theta=-RJT\)。在 \(\phi=0,2\pi\) 两端，\(V\) 和 \(T\) 都相同，故边界项 \([\langle V,T\rangle]_0^{2\pi}\) 相消。由\eqref{g1:eq:length-first-variation}，
	\[
		\frac{\mathrm{d}L}{\mathrm{d}\theta}
		=-\int_{\gamma_\theta}k_g\langle V,JT\rangle\,\mathrm{d}\ell
		=\frac{\cos\theta}{\sin\theta}\,(2\pi R\sin\theta)
		=2\pi R\cos\theta.
	\]
	这与直接对 \(L=2\pi R\sin\theta\) 求导一致。北半球 \(\theta<\pi/2\) 中，增大 \(\theta\) 是向纬线的右侧移动，长度增加；符号同时核对了 \(JT\) 的方向。
\end{solution}

\begin{exercise}
	\label{g1:ex:revolution-geodesic-integral}
	设旋转曲面由 \(X(s,\theta)=(r(s)\cos\theta,r(s)\sin\theta,z(s))\) 给出，其中 \(r>0\)、\(r'^2+z'^2=1\)。证明其仿射参数测地线满足 \(r(s)^2\dot\theta=\text{常数}\)。判定参数曲线 \(\theta=\theta_0\) 和 \(s=s_0\) 何时在适当参数下为测地线，并解释守恒量与旋转对称的关系。
\end{exercise}
\begin{solution}
	度量为 \(\mathrm{d}s^2+r(s)^2\mathrm{d}\theta^2\)。由\eqref{g1:eq:christoffel-metric}，非零的 Christoffel 系数至多为
	\[
		\Gamma^s_{\theta\theta}=-rr',\qquad
		\Gamma^\theta_{s\theta}=\Gamma^\theta_{\theta s}=\frac{r'}r.
	\]
	测地线方程于是为
	\[
		\ddot s-rr'\dot\theta^2=0,\qquad
		\ddot\theta+2\frac{r'}r\dot s\dot\theta=0.
	\]
	第二式乘 \(r^2\)，正是 \(\frac{\mathrm{d}}{\mathrm{d}t}(r(s)^2\dot\theta)=0\)。当 \(\theta=\theta_0\) 时，只需令 \(s=at+b\)，所以每条子午线都是测地线轨迹。对于纬线 \(s=s_0\)，第一式成为 \(r(s_0)r'(s_0)\dot\theta^2=0\)。非常值纬线满足 \(\dot\theta\ne0\)，故它为测地线当且仅当 \(r'(s_0)=0\)；此时第二式要求 \(\theta\) 匀速变化。

	旋转把 \(\theta\) 变为 \(\theta+\alpha\)，而度量系数完全不变。守恒量可写成 \(\langle\dot\gamma,\partial_\theta\rangle=r^2\dot\theta\)，记录速度在旋转方向上的内积分量。这里的守恒不是仅凭对称直觉宣布的：上面对测地线方程的求导已经直接证明它。若曲线为单位速，且与子午线方向的有向角为 \(\beta\)，则 \(\sin\beta=r\dot\theta\)，守恒式也写为 \(r\sin\beta=\text{常数}\)。
\end{solution}

\begin{exercise}
	\label{g1:ex:length-energy-parametrization}
	在 Euclid 平面中取 \(\gamma(t)=(t+t^2,0)\)、\(0\le t\le1\)。证明它在固定端点下使长度驻定，但不使能量驻定；给出一个固定端点变分，使能量的一阶导数不为零。再证明在连接 \((0,0)\)、\((2,0)\) 且定义于 \([0,1]\) 的分片 \(C^1\) 曲线中，能量最小值为 \(2\)，并求出达到最小值的曲线。
\end{exercise}
\begin{solution}
	曲线速度 \(\dot\gamma=(1+2t,0)\) 始终向右，所以单位切向量 \(T=(1,0)\) 恒定，\(\mathrm{D}_tT=0\)。固定端点时，长度第一变分的边界项和积分项都为零。但 \(\ddot\gamma=(2,0)\ne0\)，所以能量不驻定。更直接地，取
	\[
		F(s,t)=(t+t^2+s\,t(1-t),0),\qquad V=(t(1-t),0).
	\]
	端点保持不变，而
	\[
		\left.\frac{\mathrm{d}E}{\mathrm{d}s}\right|_{s=0}
		=-2\int_0^1t(1-t)\,\mathrm{d}t=-\frac13.
	\]
	这条曲线的长度确为 \(2\)，但能量为 \(\frac12\int_0^1(1+2t)^2\,\mathrm{d}t=13/6\)。

	对任意竞争曲线 \(\eta\)，有 \(L(\eta)\ge2\)，故\cref{g1:prop:length-energy}给出 \(E(\eta)\ge L(\eta)^2/2\ge2\)。等号要求沿两端之间的线段单向行进且等速。也可直接对向量 \(\dot\eta\) 使用积分的 Cauchy--Schwarz 不等式：
	\[
		4=\left|\int_0^1\dot\eta(t)\,\mathrm{d}t\right|^2
		\le\int_0^1|\dot\eta(t)|^2\,\mathrm{d}t=2E(\eta).
	\]
	此处等号要求 \(\dot\eta\) 几乎处处为同一向量，端点条件强制该向量是 \((2,0)\)。积分后唯一得到 \(\eta(t)=(2t,0)\)，其能量恰为 \(2\)。
\end{solution}

\begin{exercise}
	\label{g1:ex:concave-boundary-angles}
	在定向平面中，按顺序连接
	\[
		(0,0),\ (2,0),\ (2,1),\ (1,1),\ (1,2),\ (0,2),\ (0,0).
	\]
	证明这是一条正向简单闭多边形边界，求六个内角和有向外角，并用局部 Gauss–Bonnet 公式核验。若把每个外角误取为两个前进切向量之间的非负夹角，会得到什么错误结果？
\end{exercise}
\begin{solution}
	该多边形围成集合 \(D=([0,2]\times[0,1])\cup([0,1]\times[0,2])\)，各条边除相邻端点外互不相交。逐边看去，\(D\) 都在左侧，故方向为正。顶点 \((1,1)\) 处，前进方向由向左变为向上，是向右转 \(\pi/2\)；这里内部占据另外三个象限，内角为 \(3\pi/2\)，外角为 \(-\pi/2\)。其余五个顶点都向左转 \(\pi/2\)，内角与外角均为 \(\pi/2\)。

	平面 \(K=0\)，每条直线边的 \(k_g=0\)，而外角之和为 \(5\pi/2-\pi/2=2\pi\)，正好满足公式；内角之和为 \(5\pi/2+3\pi/2=4\pi=(6-2)\pi\)。若六处都只取非负夹角 \(\pi/2\)，就会算得总转角 \(3\pi\)，甚至不再是 \(2\pi\) 的整数倍。错误在于抹掉了凹角处向右转的负号。
\end{solution}

\begin{exercise}
	\label{g1:ex:spherical-cap-boundary}
	半径 \(R>0\) 的球面取外向定向。固定 \(0<\varphi_0<\pi\)，令 \(D\) 是余纬 \(0\leq\varphi\leq\varphi_0\) 的北球冠，其边界为
	\[
		\gamma(\theta)=R(\sin\varphi_0\cos\theta,\sin\varphi_0\sin\theta,\cos\varphi_0),
		\qquad 0\leq\theta\leq2\pi.
	\]
	验证参数增大的方向是球冠的正向边界，计算其测地曲率、边界积分与球冠面积，核验 Gauss–Bonnet 公式。解释 \(\varphi_0>\pi/2\) 时边界项为何为负。
\end{exercise}
\begin{solution}
	沿边界取单位切向量和子午方向单位向量
	\begin{align*}
		T&=(-\sin\theta,\cos\theta,0),\\
		e_\varphi&=(\cos\varphi_0\cos\theta,\cos\varphi_0\sin\theta,-\sin\varphi_0).
	\end{align*}
	有 \(e_\varphi\times T=n=\gamma/R\)，所以 \(JT=-e_\varphi\)，指向余纬减小的球冠内部；边界方向确为正。由 \(\mathrm{d}s=R\sin\varphi_0\,\mathrm{d}\theta\)，空间导数为
	\[
		\frac{\mathrm{d}T}{\mathrm{d}s}
		=-\frac{(\cos\theta,\sin\theta,0)}{R\sin\varphi_0}.
	\]
	与切向量 \(JT\) 作内积时，法向分量不作贡献，因此
	\[
		k_g=\left\langle\frac{\mathrm{d}T}{\mathrm{d}s},-e_\varphi\right\rangle
		=\frac{\cot\varphi_0}{R},\qquad
		\int_{\partial D}k_g\,\mathrm{d}s=2\pi\cos\varphi_0.
	\]
	球面面积元是 \(R^2\sin\varphi\,\mathrm{d}\varphi\,\mathrm{d}\theta\)，所以
	\[
		\operatorname{Area}(D)=2\pi R^2(1-\cos\varphi_0),\qquad
		\int_DK\,\mathrm{d}A=2\pi(1-\cos\varphi_0).
	\]
	两积分相加为 \(2\pi\)。球冠避开南极，包含在删去南极的球面坐标图中，且无角点，故这里适用局部公式。这个图就是\cref{g1:prop:sphere-simply-connected}证明中的球极投影：从南极投向赤道平面，半径 \(R\) 时写为 \((x,y,z)\mapsto(Rx/(R+z),Ry/(R+z))\)；那里给出的有理逆映射经缩放后适用，正反映射的分母均非零，因而也是光滑坐标图。

	当 \(\varphi_0>\pi/2\) 时，球冠超过半球，曲率面积积分大于 \(2\pi\)，故边界积分必须为负。直接计算也显示 \(\cot\varphi_0<0\)：边界加速度的切向分量指向余纬增大的一侧，即球冠外侧。负号描述的是在曲面内的弯向，与这条空间圆周的非负曲率不是同一件事。
\end{solution}

\begin{exercise}
	\label{g1:ex:nonpositive-geodesic-bigon}
	设 \(D\) 是有向光滑曲面上一张坐标图内的紧闭圆盘区域，\(K\leq0\) 在 \(D\) 上成立。
	\begin{enumerate}
		\item 若 \(\partial D\) 由 \(m\) 条测地线弧组成，各内角 \(\beta_j\in(0,2\pi)\)，证明 \(\sum_j\beta_j\leq(m-2)\pi\)，并说明何时取等号。
		\item 证明两条测地线弧不能以两个不同公共端点围成这样的圆盘，并在两个端点具有非零且非整周的内角。证明 \(D\) 的整条光滑边界也不可能是闭测地线。
	\end{enumerate}
\end{exercise}
\begin{solution}
	1. 每条测地边上 \(k_g=0\)，由\cref{g1:thm:local-gauss-bonnet}，
	\[
		\sum_{j=1}^m\beta_j=(m-2)\pi+\int_DK\,\mathrm{d}A\leq(m-2)\pi.
	\]
	等号成立当且仅当 \(\int_DK\,\mathrm{d}A=0\)。由于 \(K\) 连续且非正，若内点处 \(K<0\)，某个正面积小邻域的积分就严格为负；故等号迫使所有内点 \(K=0\)，再由连续性得到 \(K\) 在闭区域上恒为零。反之 \(K\equiv0\) 显然给出等号。

	2. 两条边的情形是 \(m=2\)，第一问会给出 \(\beta_1+\beta_2\leq0\)，与两个内角都为正矛盾。若整条边界是光滑闭测地线，则没有角点、\(k_g=0\)，公式给出 \(\int_DK\,\mathrm{d}A=2\pi\)，同样矛盾。这些结论只排除题设的圆盘边界；例如平圆柱上的闭测地线不围圆盘，并不受此排除。
\end{solution}

\begin{exercise}
	\label{g1:ex:conformal-boundary-gauss-bonnet}
	设 \(D\subset\mathbb R^2\) 是具有 \(C^2\) 正则简单闭边界的紧圆盘区域，\(\varphi\) 在其邻域光滑。在该邻域上给定度量 \(g=e^{2\varphi}(\mathrm{d}x^2+\mathrm{d}y^2)\)。记 \(s_0\)、\(t\)、\(k_0\) 分别为正向边界的 Euclid 弧长、单位切向量和有向曲率，\(\nu=Jt\) 为 Euclid 向内单位法向。
	\begin{enumerate}
		\item 用标架 \(e_1=e^{-\varphi}\partial_x\)、\(e_2=e^{-\varphi}\partial_y\) 求联络一形式，并证明
		\[
			k_g\,\mathrm{d}s_g=(k_0-\partial_\nu\varphi)\,\mathrm{d}s_0.
		\]
		\item 用 \(K\,\mathrm{d}A=-(\varphi_{xx}+\varphi_{yy})\,\mathrm{d}x\,\mathrm{d}y\) 与平面 Green 公式直接证明 \(\int_DK\,\mathrm{d}A+\int_{\partial D}k_g\,\mathrm{d}s_g=2\pi\)，说明度量改变时哪两项互相补偿。
	\end{enumerate}
\end{exercise}
\begin{solution}
	1. 将 \(g_{11}=g_{22}=e^{2\varphi}\)、\(g_{12}=0\) 代入\eqref{g1:eq:christoffel-metric}，得到
	\begin{align*}
		\nabla_{\partial_x}e_1&=-\varphi_y e_2,&
		\nabla_{\partial_y}e_1&=\varphi_x e_2.
	\end{align*}
	例如 \(\nabla_{\partial_x}\partial_x=\varphi_x\partial_x-\varphi_y\partial_y\)，对系数 \(e^{-\varphi}\) 求导后，\(\partial_x\) 分量恰相消；另一式同理。因此
	\[
		\omega=-\varphi_y\,\mathrm{d}x+\varphi_x\,\mathrm{d}y.
	\]
	写 \(t=(\cos\theta,\sin\theta)\)，则度量单位切向量为 \(T=e^{-\varphi}t=\cos\theta\,e_1+\sin\theta\,e_2\)，所以两种度量下相对于相应标架的角度相同，而 \(\mathrm{d}s_g=e^\varphi\mathrm{d}s_0\)。由\eqref{g1:eq:tangent-connection-angle}，
	\[
		k_g\,\mathrm{d}s_g=\mathrm{d}\theta+\omega(t)\,\mathrm{d}s_0.
	\]
	Euclid 曲率给出 \(\mathrm{d}\theta=k_0\,\mathrm{d}s_0\)，且
	\[
		\omega(t)=-\varphi_y\cos\theta+\varphi_x\sin\theta
		=-\partial_\nu\varphi,
	\]
	因为 \(\nu=(-\sin\theta,\cos\theta)\)。所求边界公式成立。

	2. 向外单位法向是 \(-\nu\)。将\cref{g1:lem:green-formula}用于一形式 \(-\varphi_y\,\mathrm{d}x+\varphi_x\,\mathrm{d}y\)，可得
	\[
		\int_D(\varphi_{xx}+\varphi_{yy})\,\mathrm{d}x\,\mathrm{d}y
		=-\int_{\partial D}\partial_\nu\varphi\,\mathrm{d}s_0.
	\]
	因而面积项为 \(\int_{\partial D}\partial_\nu\varphi\,\mathrm{d}s_0\)，与第一问的边界项中的法向导数部分恰好抵消。剩下
	\[
		\int_DK\,\mathrm{d}A+\int_{\partial D}k_g\,\mathrm{d}s_g
		=\int_{\partial D}k_0\,\mathrm{d}s_0=2\pi,
	\]
	最后一式是\cref{g1:thm:turning-tangent}。共形度量改变了内部的曲率面积积分，也改变了边界的测地曲率积分；边界法向导数具体说明了两者如何补偿。
\end{solution}

\begin{exercise}\label{g1:ex:gb-finite-cover}
	设 \(M\) 是紧连通光滑度量曲面，允许有分片光滑边界及有限个非零内角，\(p:\widetilde M\to M\) 是有限 \(m\) 叶光滑覆盖。给 \(\widetilde M\) 配备拉回度量 \(p^*g\)。不用三角形提升计数，利用 Gauss--Bonnet 公式证明 \(\chi(\widetilde M)=m\chi(M)\)，并说明不可定向情形为何不改变证明。
\end{exercise}
\begin{solution}
	有限叶覆盖的总空间紧，可用\cref{g1:prop:surface-finite-cover-euler}证明的第一段构造直接看出：用有限个闭包位于均匀覆盖坐标图内的小邻域覆盖 \(M\)，每个小闭包的原像是 \(m\) 个紧副本，其有限并覆盖 \(\widetilde M\)。下面只使用这一紧性结论，不使用该命题的 Euler 计数结论。

	在均匀覆盖的每一片上，\(p\) 是局部等距，故保留 \(K\)、面积测度、边界的内侧单位法向以及内角。取有限个这样的图，在它们内把积分区域分成互不重叠的可积小块；每块有 \(m\) 个等距副本，所以
	\[
		\int_{\widetilde M}\widetilde K\,\mathrm{d}\widetilde A=m\int_MK\,\mathrm{d}A,
		\qquad
		\int_{\partial\widetilde M}\widetilde k_g\,\mathrm{d}\widetilde s=m\int_{\partial M}k_g\,\mathrm{d}s.
	\]
	每个角点也恰有 \(m\) 个原像，且其外角相同，故外角和同样乘以 \(m\)。两边分别应用\cref{g1:thm:global-gauss-bonnet}，得到 \(2\pi\chi(\widetilde M)=2\pi m\chi(M)\)。面积测度、曲率和内侧测地曲率都不要求全局定向，局部等距也可以反转定向，所以证明同时适用于不可定向曲面。
\end{solution}

\begin{exercise}\label{g1:ex:gb-cutting}
	设 \(M\) 是紧光滑度量曲面，边界光滑。对下列两种切割，用 Gauss--Bonnet 公式证明相应的 Euler 示性数变化。
	\begin{enumerate}
		\item 沿内部一条光滑嵌入闭曲线 \(C\) 切开，所得曲面为 \(M_C\)。证明 \(\chi(M_C)=\chi(M)\)，并同时处理 \(C\) 为单侧曲线的情形。
		\item 沿一条光滑嵌入弧 \(A\) 切开，假设其两个不同端点在 \(\partial M\) 上，内部在 \(M^\circ\) 中，且两端都与边界横交。所得带角点曲面为 \(M_A\)。证明 \(\chi(M_A)=\chi(M)+1\)。
	\end{enumerate}
\end{exercise}
\begin{solution}
	1. 切开后每个 \(C\) 的小弧都有两个边界副本，二者的内侧单位法向互为相反数。协变加速度相同，故两侧的 \(k_g\) 相反，积分抵消。若 \(C\) 为双侧曲线，这两个副本分别接成两条边界；若为单侧曲线，它们接成一条二叶覆盖 \(C\) 的边界。后一情形仍可先沿 \(C\) 分成有限小弧，再将两侧积分配对，总和依然为零。切割不改变面积积分，也不产生角点，原有边界的积分保持不变。比较 \(M_C\) 与 \(M\) 的 Gauss--Bonnet 公式，即得 \(\chi(M_C)=\chi(M)\)。

	2. 弧内部的两侧积分仍然抵消，面积积分和原有边界各段也不变。在弧的一个端点处，原来的内角为 \(\pi\)，切割将其分为 \(\theta\) 与 \(\pi-\theta\)，其中横交保证 \(0<\theta<\pi\)。这两个新角的外角和是
	\[
		(\pi-\theta)+\bigl(\pi-(\pi-\theta)\bigr)=\pi.
	\]
	两个端点共增加 \(2\pi\)。因此 \(M_A\) 的 Gauss--Bonnet 左端比 \(M\) 增加 \(2\pi\)，右端也增加 \(2\pi\)，故 \(\chi(M_A)=\chi(M)+1\)。无论切割后是否连通，这个计算均成立。
\end{solution}

\begin{exercise}\label{g1:ex:gb-cone-points}
	设 \(M\) 是紧拓扑曲面，除有限个内部点 \(q_1,\ldots,q_r\) 外具有光滑度量，边界分片光滑且与这些点不交。假设每个 \(q_i\) 的穿孔邻域等距于
	\[
		0<\rho<\rho_i,\qquad \theta\in\mathbb R/2\pi\mathbb Z,
		\qquad g=\mathrm{d}\rho^2+a_i^2\rho^2\,\mathrm{d}\theta^2,
	\]
	其中 \(a_i>0\)。称 \(\Theta_i=2\pi a_i\) 为该点的锥角。证明
	\[
		\int_{M\setminus\{q_1,\ldots,q_r\}}K\,\mathrm{d}A
		+\int_{\partial M}k_g\,\mathrm{d}s+\sum_{v\in\partial M}\alpha_v
		+\sum_{i=1}^r(2\pi-\Theta_i)=2\pi\chi(M).
	\]
	其中边界角点仍假定内角属于 \((0,2\pi)\)。
\end{exercise}
\begin{solution}
	在每个锥点附近，局部改用角变量 \(\varphi=a_i\theta\) 就把度量写成平面极坐标度量，所以穿孔邻域的 \(K=0\)。去掉互不相交的小开圆盘 \(\rho<\varepsilon\)，得到光滑度量曲面 \(M_\varepsilon\)。由\cref{g1:prop:surface-euler-formulas}中封盘增加 \(1\) 的结论，\(\chi(M_\varepsilon)=\chi(M)-r\)。

	新边界 \(\rho=\varepsilon\) 的单位切向量可取 \(T=(a_i\varepsilon)^{-1}\partial_\theta\)。由度量系数直接计算 \(\Gamma^\rho_{\theta\theta}=-a_i^2\rho\)，所以
	\[
		\nabla_TT=-\frac1\varepsilon\partial_\rho.
	\]
	\(M_\varepsilon\) 的内侧在 \(\rho\) 增大的方向，故 \(\nu=\partial_\rho\)，\(k_g=-1/\varepsilon\)。圆周长度为 \(2\pi a_i\varepsilon\)，其边界积分就是 \(-2\pi a_i=-\Theta_i\)。这些新边界没有角点；面积曲率积分在删去的穿孔邻域内为零，因此它不依赖足够小的 \(\varepsilon\)。对 \(M_\varepsilon\) 应用\cref{g1:thm:global-gauss-bonnet}，得到
	\[
		\int K\,\mathrm{d}A+\int_{\partial M}k_g\,\mathrm{d}s
		+\sum_{v\in\partial M}\alpha_v-\sum_{i=1}^r\Theta_i
		=2\pi\bigl(\chi(M)-r\bigr).
	\]
	移项即为所求。这里的 \(2\pi-\Theta_i\) 是集中在锥点处的角亏；该点不必具有通常意义下的光滑曲率函数。
\end{solution}

\begin{exercise}\label{g1:ex:negative-curvature-area-bounds}
	设 \(M\) 为亏格 \(g\ge2\) 的紧连通可定向闭度量曲面，且 \(-b\le K\le-a<0\)，其中 \(0<a\le b\)。证明
	\[
		\frac{4\pi(g-1)}b\le\operatorname{Area}(M)\le\frac{4\pi(g-1)}a,
	\]
	并给出两端分别取等号的充要条件。
\end{exercise}
\begin{solution}
	置 \(A=\operatorname{Area}(M)\)。对 \(a\le-K\le b\) 积分，并用\eqref{g1:eq:total-curvature-genus}，有
	\[
		aA\le-\int_MK\,\mathrm{d}A=4\pi(g-1)\le bA.
	\]
	除以正数 \(a,b\) 即得面积界。下界取等号当且仅当 \(\int_M(b+K)\,\mathrm{d}A=0\)。被积函数连续且非负，若某点为正则其小邻域积分为正，所以等价于 \(K\equiv-b\)。同理，上界取等号当且仅当 \(K\equiv-a\)。这里的等号结论是针对给定曲面度量的判定，不另外断言任意面积都能实现。
\end{solution}

\begin{exercise}\label{g1:ex:degenerate-vector-indices}
	把平面的点和向量分别用复数 \(z=x+\mathrm{i}y\) 表示；这只是实数对的简写。
	\begin{enumerate}
		\item 对整数 \(m\ge1\)，求向量场 \(V(z)=z^m\) 与 \(W(z)=\overline z^{,m}\) 在原点的指数。
		\item 把球面写成两张复坐标图，重叠处用 \(w=1/z\) 换坐标。证明在 \(z\) 图上给定常速度 \(\dot z=1\)，可光滑延拓为球面上的切向量场，恰有一个零点，并求其指数。
	\end{enumerate}
\end{exercise}
\begin{solution}
	1. 沿正向圆周 \(z=re^{\mathrm{i}\theta}\)，两个单位化向量分别为 \(e^{\mathrm{i}m\theta}\) 与 \(e^{-\mathrm{i}m\theta}\)，故指数为 \(m\) 与 \(-m\)。这也表明退化零点的指数并不限于 \(\pm1\)。

	2. 先说明这套坐标确实描述标准光滑球面。可取 \(z=(X+\mathrm{i}Y)/(1-Z)\) 为去掉北极的立体投影坐标，另一图取 \(w=(X-\mathrm{i}Y)/(1+Z)\)；在单位球面重叠处，由 \(X^2+Y^2=1-Z^2\) 得 \(zw=1\)。两图的坐标函数和逆映射都光滑。

	切向量按坐标映射的微分换分量，所以在另一张图上
	\[
		\dot w=-z^{-2}\dot z=-w^2.
	\]
	右端是 \(\Re w,\Im w\) 的二次多项式，在 \(w=0\) 光滑且为零；在其余位置非零。因此所得球面场恰有北极一个零点。围绕 \(w=0\) 一周，\(-w^2\) 的角度增加 \(4\pi\)，指数为 \(2\)。使用坐标基而非正交标架不改变此计算：把坐标基正交化所产生的正对角上三角换基矩阵可经 \((1-t)A+tI\) 在可逆矩阵中变到恒等，因而单位化向量的绕数不变。总指数为 \(2=\chi(S^2)\)，与\cref{g1:thm:surface-poincare-hopf}相符；该定理并不要求球面场至少有两个零点。
\end{solution}

\begin{exercise}\label{g1:ex:conformal-torus-sign}
	设 \(f(u,v)\) 为两个变量均以 \(2\pi\) 为周期的光滑实函数。在环面上赋予度量 \(\widetilde g=e^{2f}(\mathrm{d}u^2+\mathrm{d}v^2)\)。证明：若其 Gauss 曲率在全曲面上不变号，则 \(f\) 为常数。因此任何非常数 \(f\) 都使新度量的曲率同时取正值和负值。
\end{exercise}
\begin{solution}
	环面 \(\chi=0\)，由\cref{g1:prop:curvature-sign-topology}，不变号的曲率必恒为零。\eqref{g1:eq:conformal-curvature-change}又给出
	\[
		0=\widetilde K e^{2f}=-(f_{uu}+f_{vv}).
	\]
	在周期矩形 \([0,2\pi]^2\) 上乘以 \(f\) 并分部积分；相对两边的函数值及其一阶导数相等，边界项成对抵消，故
	\[
		0=-\int_0^{2\pi}\int_0^{2\pi}f(f_{uu}+f_{vv})\,\mathrm{d}u\,\mathrm{d}v
		=\int_0^{2\pi}\int_0^{2\pi}(f_u^2+f_v^2)\,\mathrm{d}u\,\mathrm{d}v.
	\]
	连续非负函数的积分为零，使 \(f_u=f_v=0\)，所以 \(f\) 在连通平面上为常数。反之常数 \(f\) 只把平坦度量按常数缩放，曲率仍为零。若 \(f\) 非常数，曲率便不能非负，也不能非正，因而必兼有正负值。
\end{solution}
